% =====================================================================
%  The Continuum Limit as a Consequence, Not a Route
%  4D pure SU(N) Yang--Mills : one day's structural findings (2026-11-25)
%  v4 (English) — supersedes v2 / v3 for everything about the continuum limit
%  Build:  pdflatex -interaction=nonstopmode main.tex   (run twice)
%  House rule: no QED, no unlabelled assertion, no number without a source file.
% =====================================================================
\documentclass[11pt,a4paper]{article}

\usepackage[T1]{fontenc}
\usepackage[utf8]{inputenc}
\usepackage{lmodern}
\usepackage[margin=2.4cm,top=2.2cm,bottom=2.3cm,headsep=14pt]{geometry}
\usepackage{amsmath,amssymb,mathtools,bm}
\usepackage{microtype}
\usepackage[table]{xcolor}
\usepackage{booktabs,array,tabularx,longtable}
\usepackage{tcolorbox}
\tcbuselibrary{skins,breakable}
\usepackage{enumitem}
\usepackage{caption}
\usepackage{titlesec}
\usepackage{fancyhdr}
\usepackage[colorlinks=true,linkcolor=blue!45!black,urlcolor=blue!45!black]{hyperref}

% ---------- colours ----------
\definecolor{cT}{HTML}{1D7A3A}
\definecolor{cC}{HTML}{1F5FA8}
\definecolor{cM}{HTML}{B45F06}
\definecolor{cN}{HTML}{A61B1B}
\definecolor{cE}{HTML}{6A3D9A}

% ---------- label chips : copied verbatim from docs/paper-v2/main.tex:30--42 ----------
\newcommand{\labb}[2]{\fcolorbox{#1!35}{#1!8}{\textcolor{#1}{\textbf{\scriptsize #2}}}}
\newcommand{\LT}{\labb{cT}{T}}
\newcommand{\LC}{\labb{cC}{C}}
\newcommand{\LM}{\labb{cM}{M}}
\newcommand{\LN}{\labb{cN}{N}}
\newcommand{\LE}{\labb{cE}{E}}

% ---------- code : detokenize so that underscores never explode ----------
\newcommand{\brk}{\hspace{0pt}}
\newcommand{\code}[1]{\texttt{\footnotesize\detokenize{#1}}}

% ---------- source line ----------
\newcommand{\src}[1]{\par\vspace{1pt}{\footnotesize\color{black!55}\raggedright
  \textbf{Source.} #1\par}\vspace{3pt}}

% ---------- key statement ----------
\newtcolorbox{keybox}{enhanced, breakable, colback=black!3, colframe=black!45,
  boxrule=0pt, leftrule=2.4pt, arc=2pt, left=9pt, right=8pt, top=5pt, bottom=5pt,
  before skip=6pt, after skip=6pt, fontupper=\small, parbox=false}

% ---------- status box ----------
\newtcolorbox{statusbox}{enhanced, breakable, colback=black!3.5, colframe=black!30,
  boxrule=0.5pt, arc=3pt, left=10pt, right=10pt, top=6pt, bottom=6pt,
  fontupper=\footnotesize, parbox=false}

% ---------- errata helpers ----------
\newcommand{\er}[1]{\textcolor{cN}{#1}}
\newcommand{\co}[1]{\textcolor{cT}{#1}}

% ---------- maths helpers ----------
\newcommand{\Gpp}{G''(0)}
\newcommand{\Wf}{W_{\mathrm{full}}}
\newcommand{\Wg}{W_{\mathrm{geo}}}
\newcommand{\dd}{\mathrm{d}}
\newcommand{\Lm}{\Lambda_{\mathrm{MS}}}
\newcommand{\tH}{\lambda_{tH}}
\newcommand{\Var}{\operatorname{Var}}

\captionsetup{font=small,labelfont=bf,skip=5pt}
\setlist[itemize]{leftmargin=1.4em,itemsep=1pt,topsep=3pt}
\setlist[enumerate]{leftmargin=1.7em,itemsep=1pt,topsep=3pt}
\allowdisplaybreaks
\setlength{\emergencystretch}{8em}
\tolerance=3000
\setlength{\parskip}{1pt}
\setlength{\parindent}{0pt}
\setlength{\parskip}{2pt}

\titleformat{\section}{\large\bfseries}{\thesection}{0.6em}{}
\titleformat{\subsection}{\normalsize\bfseries}{\thesubsection}{0.55em}{}
\titlespacing*{\section}{0pt}{1.15em}{0.5em}
\titlespacing*{\subsection}{0pt}{0.85em}{0.35em}

\pagestyle{fancy}
\fancyhf{}
\renewcommand{\headrulewidth}{0.4pt}
\fancyhead[L]{\footnotesize\color{black!55}The continuum limit as a consequence, not a route (v4)}
\fancyhead[R]{\footnotesize\color{black!55}\thepage}
\setlength{\headheight}{13pt}

\hypersetup{pdftitle={The Continuum Limit as a Consequence, Not a Route},
            pdfauthor={Yuxuan Jiang}}

\begin{document}

% =====================================================================
%  TITLE
% =====================================================================
\begin{center}
  \vspace*{0.3cm}
  {\LARGE\bfseries The Continuum Limit as a Consequence,\\[3pt] Not a Route\par}
  \vspace{0.45em}
  {\large\bfseries everything about shape and scaling is posterior to $m>0$\par}
  \vspace{0.7em}
  {\normalsize 4D pure $\mathrm{SU}(N)$ Yang--Mills: one day's \emph{structural} findings\par}
  \vspace{0.5em}
  \rule{0.55\textwidth}{0.5pt}
  \vspace{0.55em}
  {\normalsize Yuxuan Jiang\par}
  \vspace{0.3em}
  {\small\color{black!60} Fourth edition (v4, English) \textperiodcentered{} 25 November 2026
   \textperiodcentered{} supersedes v2/v3 on everything concerning the continuum limit\par}
\end{center}

\vspace{0.2em}

\begin{statusbox}
\textbf{What this document is.} It is not a proof of the Clay conjecture; progress on Clay is
$0\%$. It reports one day's (2026-11-25) \emph{structural} findings: one identity
(\S\ref{sec:scalar}), one exact counterexample together with an exact \emph{success} inside the
scaling region (\S\ref{sec:failure}), two closures (\S\ref{sec:tension}, \S\ref{sec:loop}), one
compression of the target list (\S\ref{sec:five}), one literature verdict (\S\ref{sec:lit}) and one
errata table (\S\ref{sec:errata}). \textbf{No new species is proposed} and no step is claimed to
have been completed.

\smallskip
\textbf{Labels.} Every substantive claim carries one:
\LT{} \emph{theorem / identity} (any imported standard result is named);
\LC{} \emph{conditional theorem} (a proved implication from explicitly named hypotheses);
\LN{} \emph{no-go} (the corpus tag \code{[NG]}: a theorem or an exact counterexample showing that a
whole class of approaches cannot work);
\LM{} \emph{exact missing step} (stated as an inequality or an identity, with no slack hidden);
\LE{} \emph{numerical evidence} (measured, not proved).

\smallskip
\textbf{Provenance discipline.} Every number is followed by the file it was taken from. Where the
corpus contradicts itself, both readings are displayed and it is said which side is an identity and
which is a measurement; nothing is silently repaired.
\end{statusbox}

\vspace{0.2em}

% =====================================================================
% =====================================================================
%  §1  Abstract and positioning
% =====================================================================
\section{Abstract and positioning: the continuum-limit coordinate is a consequence, not an entrance}
\label{sec:position}

\begin{keybox}
\textbf{One sentence.} \emph{All ``shape'' and ``scaling'' arguments are logically posterior to
$m>0$; therefore the continuum-limit coordinate is not where the gap is proved.} \LN
\end{keybox}
\src{\code{LOOP-RUNBOOK.md} \S10.21(5) (verbatim: ``any scaling statement about $m$ is logically
posterior to $m>0$''), \S10.22 (``the continuum limit is not where one proves the gap; it is where
the gap's consequences live''); immune family \textbf{47} (\code{LOOP-RUNBOOK.md} line 1126); the
third admission question for a new species (ibid.\ line 1141); \code{H1B-FIVE.md} \S4.3. The corpus
tag for this verdict is \code{[NG]}.}

\begin{keybox}
\textbf{Scope, which no section below retracts.} These sections are about $\Gpp$ and about the coordinate
$m/\Lm$; \emph{they do not imply a mass gap.} Two \LT{}-level obstructions are recorded in
\code{docs/paper-v5} \S2--\S3: (i) $\Gpp=\Delta_{\mathrm{dyn}}+\Wf$ with $\Wf\ge N/12>0$, so $\Gpp>0$ is
strictly weaker than $\Delta_{\mathrm{dyn}}>0$; (ii) the harmonic-oscillator~$\otimes$~free-mode
counterexample has $\Gpp=2\omega>0$ while $m=0$. The correct target is $f_{1}^{2}$ (the
$O=\frac14\operatorname{tr}F^{2}$ pole residue) together with $(CP)$; see \code{paper-v5} \S7. \LN
\end{keybox}
\src{\code{docs/paper-v5} \S2--\S3 and \S7; \code{T01-IMPLY.md}, \code{T03-TOY.md}. Every label in the
sections below stands as stated; withdrawn is only the reading of $\Gpp>0$ as a route to $m>0$.}

\subsection{What is refuted is the ``route'' reading, not the shape}

On 2026-11-25 this campaign fixed its single objective as
\begin{equation}
  \frac{m}{\Lm} \;\ge\; c(N)\;>\;0 ,
  \qquad c \ \text{a pure number},
  \label{eq:target}
\end{equation}
and reduced it to a slope identity along $\beta$,
$\dd\ln(ma)/\dd\beta = -12\pi^{2}/(11N^{2})$ (\code{CONTINUUM-PLAN.md} \S1). The reduction itself is
an identity: $1/(2b_0g^{2}) = 12\pi^{2}\beta/(11N^{2})$ with $b_0=11N/(48\pi^{2})$. \LT
\src{\code{CONTINUUM-PLAN.md} \S1 (including the $\mathrm{SU}(2)$ coefficient $3\pi^{2}/11 = 2.691710$
and the $\mathrm{SU}(3)$ coefficient $4\pi^{2}/33 = 1.196316$); independently rechecked in
\code{LOOP-RUNBOOK.md} \S10.19 and \S10.21(1).}

But reading it as a \emph{route} (``prove the slope identity $\Rightarrow$ obtain the gap'') died the
same day, for three independent reasons. \LN
\src{\code{C1-SLOPE-EQUIV.md}, as reported in \code{LOOP-RUNBOOK.md} \S10.21; the three reasons are
listed below one by one.}

\begin{enumerate}[label=(\Alph*)]
\item \textbf{Same order $=$ tautology.} Write $\hat m := m a$. Since
      $\hat m = (m/\Lambda)(a\Lambda)$, one has $\hat m/(a\Lambda) \equiv m/\Lambda$ ---
      \emph{zero physical content}. \LT
\item \textbf{Mixed order $=$ false, in both directions.} A perfectly physical $m/\Lambda = 1$ gives
      the slope $-1.126068 \neq -1.196316$; conversely, forcing the slope to be the constant forces
      $m/\Lambda \propto \beta^{-0.4215}$, which is false as well. \LT
\item \textbf{Presupposition.} The slope identity presupposes $ma>0$; every scaling statement about
      $m$ is logically posterior to $m>0$, hence it cannot structurally be used to prove a gap. \LN
\end{enumerate}
\src{\code{LOOP-RUNBOOK.md} \S10.21(2)(3)(4)(5); the numbers $1.126068$ and $0.4215$ are the two
counterexamples of item (4) there.}

\subsection{What survives is uniqueness: the shape is a theorem, and a theorem is not a route}

The route reading is dead; the \emph{shape} is promoted to a theorem. The list of scale sources is
closed: $\Lambda$ is the only independent scale ($a$ is only a unit and enters only through
$\Lambda a$; $L$ enters only through $\Lambda L$; $\theta$ is an angle with $[\,\theta\,]=0$; $N$ is
discrete; $T$ disappears at $T=0$), and $\theta=0$ holds \emph{by definition} from the positivity of
the Wilson measure, not by hypothesis. Hence
\begin{equation}
  m \;=\; C(N,\ \text{scheme})\cdot\Lambda
  \qquad\text{is the only possible shape},
  \label{eq:unique}
\end{equation}
escape shapes $m=C/a$ and $m=C\Lambda^{2}/g$ are excluded, and $C$ may depend on $N$, on $\theta=0$
and on the scheme ($\Lm/\Lambda_I = 2.633$) but \emph{not} on $\beta,a,L,T$. \LT
\src{\code{C3-UNIQUENESS.md}, as reported in \code{LOOP-RUNBOOK.md} \S10.19 and in the first row of
the table in \S10.22; the scheme ratio $2.633$ is from \code{L2-GLUEBALL-LAMBDA.md}.}

\noindent
This is exactly what the title means: \eqref{eq:unique} is a theorem about how $m$ must look
\emph{given that $m$ exists}. \textbf{A theorem about shape can never serve as the route that
produces $m>0$} --- it is a consequence. \LT
\src{Ibid.; see also the third admission question in \code{LOOP-RUNBOOK.md} line 1141 (``does it
deliver $m>0$ or only some property of $m$?'').}

\subsection{The same day turned ``no work inside the continuum limit'' into five closures}

The remaining sections are not new conjectures; they cash out the verdict coordinate by coordinate.
The last row of the table is the only thing still alive.

\begin{table}[h]
\centering\footnotesize
\setlength{\tabcolsep}{4pt}
\renewcommand{\arraystretch}{1.15}
\begin{tabularx}{\textwidth}{@{}l l X l@{}}
\toprule
\textbf{\S} & \textbf{Object} & \textbf{Verdict (one line)} & \textbf{Label}\\
\midrule
\S\ref{sec:scalar} & scalar form of H1 & $\Var - E[H] - \Wf = -\Gpp$, hence $H1(\delta) \iff \Gpp \ge \delta$: there is no independent variance inequality & \LT\\
\S\ref{sec:failure} & solvable models & $\delta(\beta)=0.2051153-0.0110593\beta^{2}$, $\beta^{*}=4.30661$, but at $N=3$ the coefficient changes sign, so H1 \textbf{holds on the whole $\mathrm{SU}(3)$ scaling region}; with $\beta^{*}(N)\propto N^{3/2}$ one gets $\beta^{*}(3)\approx19.6\gg\beta_c(\mathrm{SU}(3))=5.6924$ and $\beta^{*}/\beta_c=1.87\to3.45$, so the exact failure sits outside the scaling region of \emph{both} groups, and the $\beta$-uniform reading is dead & \LE{}/\LN\\
\S\ref{sec:five} & five conditions & $C5\Rightarrow C2$; $C3/C4$ are filters; only $C5$'s rate uniformisation $=d5/(CP)=p=1$ survives & \LM\\
\S\ref{sec:tension} & string-tension side & $\sigma\propto 1/a^{2}$ diverges as a power; $\sigma=2g^{2}N m^{2}$ is scheme bookkeeping; $\sigma>0\not\Rightarrow m>0$ & \LN\\
\S\ref{sec:loop} & loop algebra & $\kappa_W=0$ excludes homogeneous monomials only (a weak result); the area law holds only in the degenerate sector; the non-degenerate sector is never elliptic & \LN\\
\S\ref{sec:lit} & literature & $m/\Lm = 6.3(3)$ at the continuum limit, three routes agreeing; $a$-uniform theorems: \textbf{zero} & \LE\\
\S\ref{sec:errata} & errata & 24 normalization / coordinate / convention errors of the day, each as ``wrong $\to$ right $\to$ source'' & ---\\
\bottomrule
\end{tabularx}
\caption{The structure of this document. The first five sections are closures, the sixth is the only
surviving target, the seventh is the external calibration, the eighth is the ledger.}
\label{tab:map}
\end{table}

\subsection{Conventions (read before \S\ref{sec:tension} and \S\ref{sec:errata})}

\begin{itemize}
\item Lattice $\beta := 2N/g^{2}$ (Wilson convention); $\tH := \beta N = 2N^{2}/g^{2}$.
\item SZZ's $\beta_A$ is \textbf{not} the same coordinate: $\beta_W = N^{2}\beta_A$. Mixing them
      misses by exactly $N^{2}$ ($\mathrm{SU}(3)$: $9$); this is errors 4 and 10 of the day. \LT
      \src{\code{BETA-CONVENTION.md} (read to arXiv:2204.12737v1 \S1.1 eq.~(1.2), \S1.2 eq.~(1.3));
      \code{LOOP-RUNBOOK.md} \S9.8.}
\item Every $\delta$ below is a \textbf{surfaced lower bound in lattice units} (dimensionless) and
      must not be read as a physical mass bound: $\Gpp_a = a^{2}m_{\mathrm{phys}}^{2}$, so ``for
      fixed $a$, $\Gpp_a \ge \lambda$'' is dimensionally wrong; the correct target is
      $\Gpp_a/a^{2} \ge \Delta_0^{2}$. \LT
      \src{\code{docs/paper-v2/ch03-existence.tex}:146--165 (\code{sec:dim}), as reported in
      \code{H1A-EXACT.md} \S1.4.}
\end{itemize}

% =====================================================================
%  §2  The scalar form of H1
% =====================================================================
\section{The scalar form of H1: $H1(\delta) \iff \Gpp \ge \delta$}
\label{sec:scalar}

This section does exactly one thing: it rewrites, symbol by symbol, a condition that is usually
presented as a variance inequality into a \emph{curvature} inequality. The rewrite costs nothing ---
it is a transposition, model independent. Its payoff is decisive: from then on, ``test H1'' means
``compute $\Gpp$ exactly'', with no remaining freedom.

\subsection{Four quantities, defined one by one}

\begin{table}[h]
\centering\footnotesize
\setlength{\tabcolsep}{5pt}
\renewcommand{\arraystretch}{1.15}
\begin{tabularx}{\textwidth}{@{}l X l@{}}
\toprule
\textbf{Quantity} & \textbf{Definition} & \textbf{Location in the corpus}\\
\midrule
$\Wg$ & $E[\Phi'']-\Var(\Phi')$, where $\Phi=-\log J$ and $J$ is the Haar Jacobian of the flat chart
$X=\log U$ ($\dd U = J\,\dd X$) & \code{ch05-theorems.tex}:196--206 (\code{eq:wgeo}); \code{app.tex}:795\\
$X$ & $-2\operatorname{Cov}(G,\Phi')$ with $G=\langle v,\nabla S_W\rangle$ & \code{ch05-theorems.tex}:230--239; \code{app.tex}:796\\
$H_{vv}$ & the $vv$ component of $\operatorname{Hess}V$, $V=S_W+\Phi$; the sign really lives in the
\emph{conditioned} fibre Hessian $K=\Pi_0\operatorname{Hess}V\,\Pi_0$ & \code{ch13-machine.tex}:58--82\\
$\Var$ & $\Var_\mu(\langle v,\nabla S\rangle \mid \xi=0)$, \textbf{conditioned} (not the unconditional variance) & \code{ch08-immune.tex}:13\\
\bottomrule
\end{tabularx}
\caption{The four sources of notation. $W_{\mathrm{full}} := \Wg + X$ merely bundles the last two
terms and carries no new content.}
\end{table}

The exact form of H1 in the paper is (\code{ch08-immune.tex}:11--19, display $(M)$, label
\code{eq:species})
\begin{equation}
  \Var_\mu(\langle v,\nabla S\rangle \mid 0) \;\le\; E_\mu[H_{vv}\mid 0] + \Wf - \delta ,
  \qquad \text{uniform in } (L,U') ,
  \label{eq:M}
\end{equation}
and its slack-free form is written in \code{app.tex}:179--202 (\code{app:sharp}):
$\Var_\mu(\langle v,\nabla V\rangle\mid \xi=0) \le E_\mu[H_{vv}\mid\xi=0]-a^{2}\Delta_0^{2}$,
together with the explicit $\delta := a^{2}\Delta_0^{2} + \Wf$. \textbf{The two displays are the same
inequality.} \LT
\src{\code{H1A-EXACT.md} \S1.1, quoting the three places in v2 verbatim; that file already checked
the equivalence of the two forms.}

\subsection{The identity and the transposition}

The key identity (\code{ch05-theorems.tex}:230--239 \code{eq:split}; \code{ch03-existence.tex}:129--134;
\code{ch13-machine.tex}:69--73):
\begin{equation}
  \Gpp \;=\; E[H_{vv}\mid 0] - \Var(\langle v,\nabla S\rangle\mid 0) + \Wg + X
        \;=\; E[H\mid 0] - \Var(G\mid 0) + \Wf .
  \label{eq:split}
\end{equation}
Transposing \eqref{eq:split} (pure algebra, valid in any model):
\begin{equation}
  \boxed{\;\Var - E[H] - \Wf \;=\; -\Gpp\;}
  \qquad\Longrightarrow\qquad
  \boxed{\;H1(\delta) \iff \Gpp \ge \delta\;}
  \label{eq:equiv}
\end{equation}
pointwise, mode by mode, exactly. \LT
\src{\code{H1A-EXACT.md} \S1.3 (lines 57--68). The paper itself says in \code{ch08-immune.tex}:19
that ``the display is then $G''\ge\delta$ in curvature form''; here that sentence is simply used to
the end.}

\begin{keybox}
\textbf{Corollary (the first structural gain of this document).} The sign of $\delta$ equals the
sign of the conditioned curvature $\Gpp$, and $\delta_{\max} = \Gpp$. Hence in any model, ``test H1''
means ``compute $\Gpp$ exactly'': \textbf{there is no independent variance inequality to prove or to
refute} --- all the information of H1 sits in the sign of the curvature. \LT
\end{keybox}
\src{\code{H1A-EXACT.md} \S1.3, last paragraph.}

\subsection{What the scalar test now returns: a certificate, not an estimate}
\label{sec:scalar-cert}

\S\ref{sec:failure} evaluates $\Gpp$ by quadrature and by Monte Carlo. This subsection records the one
place where the same number is evaluated by \emph{interval arithmetic}, that is, where its sign stops
being an estimate and becomes a machine-checked inequality.

\begin{keybox}
\textbf{Theorem (three certified values).} On a single $\mathrm{SU}(2)$ plaquette, $\beta=2.5$, flat
chart $|\xi|<2\pi$, slice $\xi_{3}=a=0$, $v=e_{3}$, $\|v\|=1$, with boundary holonomy $\theta$:
\begin{gather}
\mathrm{gap}(0)=0.729924128620548449,\qquad
\mathrm{gap}(\pi/2)=0.5429439223536897638,\notag\\
\mathrm{gap}(\pi)=0.1359947110419882110,
\label{eq:certpoints}
\end{gather}
each enclosed within $\pm10^{-55}$, where
$\mathrm{gap}(\theta):=S-D=\Gpp$ with $S=E[H+\Phi''\mid0]$ and $D=\Var(G+\Phi'\mid0)$. At $\theta=0$
the variance is not small but \emph{identically zero}: $s=\sin(\theta/2)=0$ makes $G$ vanish pointwise,
so $D\equiv0$ for every $\beta$. \LT
\end{keybox}
\src{\code{R10-AUDITCERT.md} \S0 and \S2 (the three values, the enclosures, and three independent
numerical roads: its own $256\times20$ Gauss--Legendre with a Cauchy remainder bound
$\le3.1\times10^{-82}$, \code{mp.quad} at $\mathrm{dps}=45$ agreeing to $10^{-30}$, and
\code{q04_weyl.moments} agreeing to $10^{-16}$); \code{Q11-CERT.md} \S2.2 (the $\theta=\pi$ value and
its intervals); \code{R03-CERT-THETA.md} \S2.2 (the first certificate away from $\theta=\pi$);
\code{docs/paper-v2/ch19-h1-status.tex}, subsection ``the certificate'', for the same three numbers
inside the ledger.}

\noindent
\textbf{Why this is a different kind of claim.} The pipeline has five ingredients, each checkable
(\code{Q11-CERT.md} \S2.1, \code{R03-CERT-THETA.md} \S2.1): Taylor coefficients of the four entire
integrands are \emph{exact rationals} (\code{Fraction}, and $\mathbb{Q}[\sqrt2]$ at $\theta=\pi/2$);
integration is term by term, $\int_{0}^{\pi}u^{k}\,\mathrm{d}u=\pi^{k+1}/(k+1)$, an identity for entire
functions; $\pi$ and $\sqrt2$ enter only as $800$-bit \code{libmp} intervals with directed rounding
(\code{'f'} down, \code{'c'} up); every operation carries an outward pad ($2^{-400}$, absolute in
\code{q11}, relative in \code{r03}); truncation is controlled by a Cauchy estimate ($K=260$ and
$K=400$). The only input outside pure arithmetic is \code{libmp}'s own rounding, and the pad covers it
by more than a hundred orders of magnitude. The evidence that rounding is really directed is the width
pattern rather than the endpoints: in \code{q11} the enclosure of $Z$ is $2.018\times10^{-118}$ wide and
$2.018\times10^{-118}/2^{-400}=522=2(K_{\max}+1)$, exactly the per-operation accumulation, whereas a
``high precision, roughly'' computation would give width $0$ or ulp noise. Two real bugs were found by
using $\theta=\pi$ as a regression: a power index in $\sin^{4}u/(4u^{3})$ (which corrupted only
$\Var(g)$, caught because the regression must reproduce \code{q11}'s $0.0110592937$) and an
\emph{absolute} pad that, under the catastrophic cancellation of the $\theta\neq\pi$ series, inflated
$Z$ to $\pm1.56\times10^{76}$ with the endpoints inverted --- cured by a purely relative pad, after
which the $\theta=\pi/2$ width ($3.5\times10^{-21}$) matches its own Cauchy tail
($1.77\times10^{-21}$). \LT
\src{\code{Q11-CERT.md} \S2.1--\S2.2; \code{R03-CERT-THETA.md} \S2.1--\S2.2 (both bugs, and the
regression that reproduces \code{q11} to $4\times10^{-41}$); \code{R10-AUDITCERT.md} \S2 (an
independently written interval class and a true remainder bound).}

\noindent
\textbf{From three points to the interval, and the honest grading.} Three points are not an interval.
Write $c=\cos(\theta/2)$ and $x=\cos u$; then $\sin^{2}u\,u^{-1}\mathrm{d}u=w(x)\,\mathrm{d}x$, so the
slice measure is a one-parameter exponential family,
\begin{equation}
\nu_{c}\propto e^{\beta cx}w(x)\,\mathrm{d}x,\qquad x\in(-1,1),
\qquad\Longrightarrow\qquad
\partial_{c}E_{\nu}[f]=\beta\,\mathrm{Cov}_{\nu}(f,x),
\label{eq:expfam}
\end{equation}
and $\partial_{c}\mathrm{gap}$ splits into five terms: $T_{1}=\tfrac{\beta}{2}E[g]>0$;
$T_{2}=\tfrac{\beta^{2}c}{2}\mathrm{Cov}(g,x)>0$ (rearrangement: $g$ and $x$ are both strictly
decreasing in $u$); $T_{3}=\beta\,\mathrm{Cov}(\Phi'',x)<0$, the only negative term, negative because
$\Phi''$ is strictly increasing on $(0,\pi)$; $T_{4}=2\beta^{2}c\Var(g)>0$; and
$T_{5}=-\beta^{3}(1-c^{2})[\mathrm{Cov}(g^{2},x)-2E[g]\mathrm{Cov}(g,x)]$, \emph{whose sign is open}.
On $501$ interior points of $c\in(0,1)$, $\partial_{c}\mathrm{gap}\in(0.44361829,0.63898109)$ with zero
sign flips and no interior extremum, so the gap is strictly decreasing in $\theta$ and its unique worst
point is $\theta=\pi$, where \eqref{eq:certpoints} certifies $+0.1360$. \textbf{The grading is therefore
split, and the split is the point}: the three values are \LT; the statement for all
$\theta\in[0,\pi]$ is \LC, conditional on the numerical sign of $T_{5}$. \LC
\src{\code{R06-GAPMONO.md} \S1 (the exponential-family rewriting and the covariance formula), \S2 (the
five terms and their numerical table), \S3.1--\S3.2 ($\Phi''$ strictly increasing, then the $501$-point
sign measurement and the $2.9\times10^{-11}$ agreement with a central difference). \textbf{The
complement is cheap and explicit}: a grid of certified points at rational $c\in\{0,\tfrac15,\dots,1\}$
together with one bound on $|\mathrm{gap}''|\le0.597$ covers the interval as soon as
$L\,\dd c<0.13599$, which holds from $K=6$ ($0.1148$) and fails at $K=5$ ($0.1397$); four more points
therefore make all of $[0,\pi]$ unconditional \LT{} without any monotonicity (\code{R06-GAPMONO.md}
\S6).}

\noindent
\textbf{And the certificate is a window, not a theorem about the theory.} At $\theta=\pi$ one has
$S=1/(4Z)$, independent of $\beta$, while $D=\beta^{2}V_{0}$ with $V_{0}=0.0110592937729736$, so the
same computation gives a positive sign exactly for $\beta<\beta^{*}=4.30660785\ldots$: $\beta=2.5$ has
a $41.9\%$ margin and $\beta=4.4$ gives $-0.00899$. The certificate and the exact counterexample of
\S\ref{sec:failure} are one and the same computation at two couplings. No number above is scaled to
$a\to0$, and none of them is a statement about four-dimensional Yang--Mills. \LN
\src{\code{Q11-CERT.md} \S5 (the $\beta$ table and the margin); \eqref{eq:delta-beta} of
\S\ref{sec:failure} for the identical closed form.}

\subsection{The eight strict lemmas the scalar test rests on}
\label{sec:scalar-lemmas}

The certificate is not an isolated computation. Eight statements are proved rather than measured, and
each is useful only together with the scope attached to it in Table~\ref{tab:lemmas}.

\begin{table}[h]
\centering\footnotesize
\setlength{\tabcolsep}{4pt}
\renewcommand{\arraystretch}{1.1}
\begin{tabularx}{\textwidth}{@{}>{\raggedright\arraybackslash}p{0.35cm}
  >{\raggedright\arraybackslash}p{0.12cm} X >{\raggedright\arraybackslash}p{3.5cm}@{}}
\toprule
\textbf{\#} & & \textbf{Statement} & \textbf{Scope / source}\\
\midrule
1 & & $\Phi''(X;v)\ge\tfrac{N}{12}\langle v,v\rangle$ for every $X$ in the flat chart and \emph{every}
direction; sharp, attained only at $X=0$ & all $N$, all measures, all $\beta$;
\code{P16-PHIQUANT.md} \S3--\S3.1\\
2 & & $E[H_{vv}\mid0]\ge0$, indeed $H_{vv}=\tfrac{\beta c}{2}g(u)\ge0$ pointwise & $\theta\le\pi$ only;
false for $\theta>\pi$; \code{P05-CINF.md} \S1.2\\
3 & & $E[H_{vv}]=E[G_{v}^{2}]$, hence $E[H]-\Var(G)=E[G]^{2}\ge0$ & measure $e^{-S_{W}}\mathrm{dHaar}$
(no $\Phi$); \code{Q15-SPECTRAL.md} \S4\\
4 & & $H=\tfrac12\cot(\theta/2)G$ pointwise & $a=0$ slice only; \code{Q04-WEYL.md} (2.2)\\
5 & & $\Var(w)\le\tfrac12\langle w\rangle$, i.e.\ $2\Var\le\langle w\rangle$ & unconditional in
$(\beta,c)$, three independent proofs; \code{R02-CRIT.md} \S2, \code{R05-WCF.md} \S6.1,
\code{R09-FLUCT.md} \S4\\
6 & & $\langle w\rangle-2\Var(w)\ge2\langle w\rangle^{2}\ge0$ & unconditional, strictly stronger than
5; \code{R05-WCF.md} \S6.1\\
7 & & $\Gpp=m_{G'}+E[\Phi'']-\Var(\Phi')-\mathrm{Cov}(G,\Phi')$ & algebra only; the skeleton every
number is measured against; \code{P21-REBUILD.md} \S0\\
8 & & $H=K_{0}G+\beta\Delta$, $K_{0}=\tfrac12\cot(\theta/2)$,
$\Delta=[t\partial_{a}^{2}-c^{2}/(4t)](aw)$ & exact for all $a$; $\Delta$ odd, first order $O(a)$;
\code{R07-AFFINE.md} (4.1)\\
\bottomrule
\end{tabularx}
\caption{The eight strict lemmas. ``Scope'' is load-bearing: a lemma quoted without it is false
somewhere in the same corpus.}
\label{tab:lemmas}
\end{table}

\noindent
Three of the eight are worth one sentence each, because each one \emph{closes} a route rather than
opening one. Lemma 1 is the curvature floor: subtracting the quadratic part of $\Phi$ leaves
$h(X)=\sum_{i<j}\varphi(\lambda_{i}-\lambda_{j})$ with $\varphi=g-d^{2}/12$ convex, hence a symmetric
convex spectral function, hence convex in $X$ (Davis 1957, checked at bibliographic level only, with an
elementary hand proof at $N=2$); because $\Phi$ is additive over links the floor holds for every
direction, which is the step earlier ledgers claimed without having. It is a floor and not a witness:
at $a=0.75$, $\beta=2.5$, $\theta=\pi$ one has $\Phi''-\Var=+0.176>0$ while $\Gpp=-0.0059<0$, the
deficit being $E[H_{vv}]=-0.1816$; and the floor cannot buy back $-\beta E[C_{\mathrm{geo}}]$
($0.167+0.08$ against $-0.947$). Lemma 3 is Stein: left invariance makes $\mathrm{div}_{\mathrm{Haar}}v=0$,
so $E[H_{vv}]=E[G_{v}^{2}]$ and the difference is a square --- and consequently the Hessian spectrum
enters only as $E[\lambda_{k}]=E[g_{k}^{2}]$ and cancels, so no spectral condition is equivalent to the
target; the one named spectral criterion, $\Lambda\preceq I$, is \emph{false} here
($\lambda_{k}=\beta\cos\varphi$ threefold degenerate, $\lambda_{\max}=2.5$, $P(\lambda<0)=0.105$) while
the target holds and is an equality. Lemma 8 is the exact replacement for the affine fit
$H=K(a)G+C(a)$ that the corpus had been using at $R^{2}>0.9992$: two Wronskians are nonzero at first
order in $a$, the difference $H-KG$ has no constant term in the basis
$\{\sin(\rho/2),\cos(\rho/2),1\}$, so a strict affinity would force a strict proportionality, which the
explicit ratio $Q_{H}/Q_{G}=K_{0}+3/a-3a/\rho^{2}$ kills. The residual is not noise: it is
$O(a)$ with $\mathrm{relres}/a\to1.43\times10^{-3}$, unchanged when the quadrature is raised fourfold.
\LT{}/~\LE
\src{\code{P16-PHIQUANT.md} \S3, \S3.1, \S6; \code{Q15-SPECTRAL.md} \S3--\S4;
\code{R07-AFFINE.md} \S3--\S5. The Davis reference is C.~Davis, \emph{All convex invariant functions of
hermitian matrices}, Arch. Math. \textbf{8} (1957) 276--278, DOI 10.1007/BF01898787, cited at
bibliographic level (the text was not read).}

\subsection{Dimensional discipline: $\delta$ is a surfaced bound, not a mass bound}

$\Gpp_a = a^{2}m_{\mathrm{phys}}^{2}$. Hence ``for fixed $a$, $\Gpp_a \ge \lambda$ with $\lambda$ a
positive constant'' is dimensionally wrong; the correct target is
$\Gpp_a/a^{2}\ge\Delta_0^{2}$. Every $\delta$ in this document is a dimensionless surfaced lower
bound in lattice units, and its $(L,a)$-uniformity is precisely the \emph{open} point; it must not be
read as a physical mass bound. \LT
\src{\code{docs/paper-v2/ch03-existence.tex}:146--165 (\code{sec:dim}); \code{H1A-EXACT.md} \S1.4.}

\subsection{What the rewrite changes}

Before the rewrite, H1 looks like a variance-estimation task: bound $\Var(G\mid 0)$ above, bound
$E[H_{vv}\mid 0]$ below, and combine into a positive $\delta$. After the rewrite there is only one
exit left: \emph{prove that the conditioning point lies in a vacuum with $\Gpp>0$}, and control the
Jacobian term $\Phi''$ (that is, the sign of $\Wg$).

This is not a matter of taste. \S\ref{sec:failure} exhibits a 3-dimensional model whose variance
term grows like $\beta^{2}$ while the Haar term $\Wg$ is $\beta$-independent, the two cancelling
exactly at $\beta^{*}=4.30661$; and it exhibits the opposite phenomenon at $N=3$, where the same
coefficient flips sign and H1 holds for every $\beta>0$, that is, on the whole $\mathrm{SU}(3)$ scaling
region. Any argument that only estimates the
variance and never asks whether the conditioning point is a minimum will get the sign wrong on one
side or the other. \LN
\src{\code{H1A-EXACT.md} \S2.3 for the mechanism (Savvidy concavity); the $N=3$ sign flip is the new
measurement of \S\ref{sec:failure}.}

\noindent
One sentence closes this section: \textbf{H1 is not a new condition; it is the single-scalar form of
$\Gpp>0$} (positivity of the second derivative of the free energy, i.e.\ of the susceptibility). It
is necessary but not sufficient, because it is completely blind to a thin atom (small $f_1^{2}$). \LT
\src{\code{LOOP-RUNBOOK.md} \S11.39, first paragraph; the machine-checked version of ``necessary but
not sufficient'' is the module \code{H1Sharpening} in \code{docs/paper-v2} (abstract, item (ii)).}

% =====================================================================
%  §3  Solvable models: exact success inside the scaling region, exact failure outside it
% =====================================================================
\section{Solvable models: exact success inside the scaling region, exact failure outside it}
\label{sec:failure}

\S\ref{sec:scalar} says that the whole content of H1 is $\Gpp\ge\delta$. This section evaluates
$\Gpp$ exactly, first at $N=2$ on a 3-dimensional model, then at $N=3$ on the same model by Monte
Carlo. The two answers are different in sign, and the difference matters: \textbf{at $N=2$ there is
an exact zero, but it lies outside the window where the model is read; at $N=3$ the coefficient flips
sign, so there is no zero at all and H1 holds on the entire $\mathrm{SU}(3)$ scaling region.}

\subsection{Model 1: a single $\mathrm{SU}(2)$ plaquette, closed form}

Take $U=\exp(\xi)$, $\xi\in\mathbb{R}^{3}$, $|\xi|<2\pi$; boundary holonomy
$B=\exp(i\theta\sigma_3/2)$; action and flat-chart measure
\begin{gather}
  S=\beta\Big(1-\cos\tfrac{\rho}{2}\cos\tfrac{\theta}{2}
            +\sin\tfrac{\rho}{2}\sin\tfrac{\theta}{2}\,u\Big),
  \qquad u=\frac{\xi_3}{\rho},\notag\\
  \mathrm{d}\mu = e^{-S}J\,\dd^{3}\xi,
  \qquad J=\frac{\sin^{2}(\rho/2)}{\rho^{2}},
  \qquad \Phi=-\log J ,
  \label{eq:model1}
\end{gather}
with $v=e_3$ and conditioning $\langle v,\xi\rangle = a_0 = 0$. \LT
\src{\code{H1A-EXACT.md} \S2.1 (same convention as \code{research/v3/species/w7/na_plaquette.py}).}

On the slice there are two exact simplifications, and everything below rests on them: \LT
\begin{enumerate}[label=(\roman*)]
\item $\Phi$ is radial, so $\Phi'\equiv 0$ on the slice, hence $\Var(\Phi')=0$ and
      $X=-2\operatorname{Cov}(G,\Phi')=0$;
\item $\rho'=0$ and $u''=0$, hence $H_{vv}=\beta c\sin(r/2)/(2r)$ with $c=\cos(\theta/2)$; so at
      $\theta=\pi$ one has $c=0$ and $E[H]=0$ \emph{exactly} (not a small number).
\end{enumerate}
\src{\code{H1A-EXACT.md} \S2.1, lines 84--88.}

\subsection{$\theta=\pi$: closed form, zero, and a table}

At $\theta=\pi$ the vacuum sits at $\xi_3=-\pi$ (far from the conditioning point); on the slice
$S\equiv\beta$ independently of $r$, so the measure is \emph{exactly} Haar. Hence \LT
\begin{gather}
  \Var=\beta^{2}V_0,\qquad
  V_0=\Var_{\mathrm{Haar}}\!\Big(\tfrac{\sin(r/2)}{r}\Big)=0.0110593,\notag\\
  E[H]=0,\qquad \Wg = 0.2051153\ (\beta\text{-independent}),\qquad
  X=0 ,\notag
\end{gather}
\begin{equation}
  \boxed{\;\delta_{\max}(\beta)=\Gpp(\beta)=0.2051153-0.0110593\,\beta^{2},
  \qquad \beta^{*}=\sqrt{W_0/V_0}=4.30661\;}
  \label{eq:delta-beta}
\end{equation}
\src{\code{H1A-EXACT.md} \S2.1, lines 93--95; $W_0=1/(2\,\mathrm{Cin}(2\pi))=0.20511530$, see
\code{docs/paper-v2/ch15-uvface.tex}:185--197 (\code{eq:v3anchor}).}

\begin{table}[h]
\centering\footnotesize
\setlength{\tabcolsep}{7pt}
\renewcommand{\arraystretch}{1.1}
\begin{tabular}{@{}r r r r r@{}}
\toprule
$\beta$ & $\Var=\beta^{2}V_0$ & $E[H]$ & $\Wf=\Wg$ & $\delta_{\max}=\Var-E[H]-\Wf$\\
\midrule
$0.0$      & $0.000000$ & $0.000000$ & $0.205115$ & $+0.205115$\\
$0.5$      & $0.002765$ & $0.000000$ & $0.205115$ & $+0.202350$\\
$1.0$      & $0.011059$ & $0.000000$ & $0.205115$ & $+0.194056$\\
$\mathbf{2.5}$ & $\mathbf{0.069121}$ & $\mathbf{0.000000}$ & $\mathbf{0.205115}$ & $\mathbf{+0.135995}$\\
$4.0$      & $0.176949$ & $0.000000$ & $0.205115$ & $+0.028167$\\
$\mathbf{4.30661}$ & $\mathbf{0.205115}$ & $\mathbf{0.000000}$ & $\mathbf{0.205115}$ & $\mathbf{0.000000}$\\
$\mathbf{6.0}$ & $\mathbf{0.398135}$ & $\mathbf{0.000000}$ & $\mathbf{0.205115}$ & $\mathbf{-0.193019}$\\
$10.0$     & $1.105929$ & $0.000000$ & $0.205115$ & $-0.900814$\\
\bottomrule
\end{tabular}
\caption{Model 1 ($\theta=\pi$, $N=2$) row by row in $\beta$. The physical point $\beta=2.5$ is on the
safe side, the zero is exactly at $\beta^{*}=4.30661$, and beyond it $\delta<0$. Numbers from
\code{H1A-EXACT.md} \S2.1.}
\label{tab:model1}
\end{table}

The same model at $\theta=0$ (vacuum at the conditioning point) gives the other side: on the slice
$s=0$, so $G\equiv0$ and $\Var=0$ exactly, while $E[H]>0$; hence
$\delta=+0.290877\,(0.5)$, $+0.729924\,(2.5)$, $+1.590639\,(6.0)$, $+2.587494\,(10.0)$, the margin
increasing \emph{monotonically} in $\beta$. \LT
\src{\code{H1A-EXACT.md} \S2.1, the $\theta=0$ table (lines 110--115).}

\begin{keybox}
\textbf{First reading.} At $\theta=0$ H1 is true for all $\beta$; at $\theta=\pi$ the margin falls
like $\beta^{2}$, vanishes exactly at $\beta^{*}=4.30661$ and is negative beyond. H1 is therefore not
$\beta$-uniform: its truth value depends on whether the \emph{vacuum sits at the conditioning point},
not on any estimation technique. \LN
\end{keybox}

\subsection{Mechanism: exact Savvidy concavity}

At $\theta=\pi$ the action $S=\beta\big(1+\sin(\rho/2)u\big)$ is \emph{linear} in $u=\xi_3/\rho$, so
the conditioning point $\xi_3=0$ is not a minimum (the minimum is at $\xi_3=-\pi$). On that slice $S$
does not depend on $r$, hence the measure is Haar, hence all the bending comes from the Jacobian, and
the $\beta^{2}$ variance term drives $\Gpp$ negative. In one line: \textbf{``curvature at the
conditioning point'' $=$ ``stability of that point''}. \LT
\src{\code{H1A-EXACT.md} \S2.3.}

This is exactly the shape of the wall recorded in \code{CLAUDE.md} \S5: $\operatorname{Hess}$ of
$-D^{2}+2adF_B$ is indefinite, with W4 measuring $7.1\%$ negative modes. The solvable model gives the
\textbf{exact scalar version} of it: $\delta<0 \iff$ the vacuum is not at the conditioning point. \LT
\src{\code{H1A-EXACT.md} \S2.3; the negative-mode fraction is from \code{CLAUDE.md} \S5.}

\subsection{Four models: sign table and three verdicts}

\begin{table}[h]
\centering\footnotesize
\setlength{\tabcolsep}{5pt}
\renewcommand{\arraystretch}{1.12}
\begin{tabularx}{\textwidth}{@{}>{\raggedright}p{3.6cm} r r r r l@{}}
\toprule
\textbf{Model} & $\Var$ & $E[H]$ & $\Wf$ & $\delta$ & \textbf{Clause}\\
\midrule
1 \; $\theta=\pi,\ \beta=2.5$ & $0.069121$ & $0$ & $0.205115$ & $+0.135995$ & (a)\\
1 \; $\theta=\pi,\ \beta=6$ & $0.398135$ & $0$ & $0.205115$ & $\mathbf{-0.193019}$ & \textbf{(b)}\\
1 \; $\theta=\pi,\ \beta=\beta^{*}$ & $0.205115$ & $0$ & $0.205115$ & $0$ & (c)\\
1 \; $\theta=0,\ \beta=2.5$ & $0$ & $0.553061$ & $0.176864$ & $+0.729924$ & (a)\\
2 \; 2D YM ($t\to\infty$, $\varphi=1.5$) & $0$ & $\Wg-\delta$ & $\Wg$ & $+2.0101$ & (a)\\
3 \; P4 periodic kernel, $\beta=2.5$ & $0$ & $0$ & $0.22196$ & $+0.22196$ & (a)\\
4 \; 1D Dirichlet string & $0$ & $\pi^{2}/L^{2}$ & $0$ & $\pi^{2}/L^{2}\to0$ & \textbf{(c)}\\
\bottomrule
\end{tabularx}
\caption{Sign table over four models (\code{H1A-EXACT.md} \S3). Clauses: (a) true; (b) H1 holds with
$\delta<0$ (exact counterexample); (c) degenerate. All numbers are deterministic quadrature or
series, never sampling.}
\label{tab:four}
\end{table}

Three verdicts, one by one: \LT
\begin{enumerate}[label=(\alph*)]
\item \textbf{(a) twice, both carried by Haar.} 2D YM gives $\delta\ge2.0101$ (uniform in the area
      $t=g^{2}A/2$); the P4 periodic kernel gives $\delta=0.22196\pm0.00011$ (uniform in $L$, positive
      intercept, $q=0.849$). But both satisfy $\Var=E[H]=0$: \emph{the whole margin comes from
      $\Wg$, with no contribution from the interaction term}.
\item \textbf{(b) an exact counterexample.} In model 1 at $\theta=\pi$, for $\beta>\beta^{*}$, H1
      holds with $\delta<0$. This is the closed form \eqref{eq:delta-beta}, not a numerical accident.
\item \textbf{(c) degenerate twice.} The 1D Dirichlet string gives $\delta=\pi^{2}/L^{2}\to0$ (no
      infrared margin); model 1 vanishes exactly at $\beta^{*}$.
\end{enumerate}
\src{\code{H1A-EXACT.md} \S3 and \S4. The P4 values are quoted from
\code{docs/paper-v2/ch15-uvface.tex}:199--213 (\code{eq:v3periodic}) and
\code{P2-PERIODIC-K0.md}, \code{P4-PERIODIC-TRANSFER.md}; 2D YM is the Migdal exact solution with
$j\le120$ and a five-point difference of step $h=2\times10^{-5}$.}

\subsection{The $N=3$ measurement: the coefficient changes sign, so there is no zero}

The question left open in \S\ref{sec:scalar} was whether $\beta^{*}$ is $N$-independent. It is not.
Repeating the same single-plaquette (patch) construction at $N=3$ by Monte Carlo gives
\begin{equation}
  \delta(0)=+0.34368\pm0.0077,\qquad
  \delta(6)=+0.63393\pm0.0676,\qquad
  \delta(10)=+0.95757\pm0.1446 ,
  \label{eq:n3}
\end{equation}
i.e.\ $\delta$ \textbf{increases monotonically in $\beta$}. Fitting the same two-parameter form
$\delta(\beta)=W_0(N)-V_0(N)\beta^{2}$ gives $V_0(3)<0$; consequently
\begin{equation}
  \boxed{\;\beta^{*}(3)\ \text{does not exist}\;}
  \qquad\Longrightarrow\qquad
  H1\ \text{holds for every }\beta>0\text{ at }N=3 .
  \label{eq:n3nobeta}
\end{equation}
In particular H1 holds on the \textbf{entire $\mathrm{SU}(3)$ scaling region}, which starts at
$\beta_c(\mathrm{SU}(3))=5.6924 \ll 10$ and includes $\beta=10$, where
$\delta=+0.95757\pm0.1446$; and it holds on the $\mathrm{SU}(3)$ target window $\beta\approx6.6$--$7.6$
of \code{A6-TARGET.md}. \LE
\src{The three points and their errors are today's single-plaquette/patch Monte Carlo
(\code{research/unconditional/w7/betastar\_n34.log}, via \code{betastar\_scaling.py}), recorded in
\code{LOOP-RUNBOOK.md} \S11.46; the fit is the same ansatz as \eqref{eq:delta-beta} with $N=3$
Haar data. \textbf{Caveat, stated here and not to be dropped: this is a Monte Carlo of a single
plaquette/patch model, not of the full theory.} It establishes the $N$-dependence of the coefficient
inside the solvable class, and nothing about $\mathrm{SU}(3)$ Yang--Mills itself. The target window
$\beta\approx6.6$--$7.6$ is from \code{A6-TARGET.md} (\S10.4 of \code{LOOP-RUNBOOK.md}).}

\begin{keybox}
\textbf{Second reading, and it replaces the first one's emphasis.} \textbf{H1 is true inside the
scaling region of both groups}: at $N=3$ there is no zero at all (\eqref{eq:n3nobeta}); at $N=2$ the
zero sits at $\beta^{*}=4.30661$, which is \emph{above} $\beta_c(\mathrm{SU}(2))=2.2986$ and above the
$\mathrm{SU}(2)$ target window $\beta\approx2.9$--$3.2$. The exact counterexample of
\eqref{eq:delta-beta} therefore lives \emph{outside} the windows in which the model is read. \LE
\end{keybox}
\src{\code{LOOP-RUNBOOK.md} \S11.39 (for the $\beta^{*}$ versus $\beta_c$ comparison at $N=2$) and
\S10.4 / \code{A6-TARGET.md} (for the target windows); the $N=3$ input is \eqref{eq:n3}.}

\noindent
\textbf{What this correction withdraws.} The previous inference --- ``$\beta_c(\mathrm{SU}(3)) >
\beta^{*}$ implies H1 fails before the $\mathrm{SU}(3)$ scaling window opens, hence H1 is the wrong
target for the Clay object'' --- is \textbf{withdrawn}: it assumed $\beta^{*}$ to be $N$-independent,
and \eqref{eq:n3nobeta} shows $\beta^{*}(3)$ does not exist. \LE
\src{The withdrawn inference is stated in \code{LOOP-RUNBOOK.md} \S11.39 (lines 1879--1885, with the
conditional ``if $\beta^{*}$ is $N$-independent or grows slower than $\beta_c(N)$''); this document
writes down its antecedent as false at $N=3$. The route-level no-go survives in the weaker and true
form below.}

\begin{keybox}
\textbf{What survives as a no-go.} The \emph{$\beta$-uniform, vacuum-independent} reading of H1 is
dead: the exact zero at $\beta^{*}=4.30661$ at $N=2$ cannot be removed by any estimate, and any proof
of H1 that does not carry vacuum-selection information (control of $\Phi''$, i.e.\ of the sign of
$\Wg$) will be broken by that 3-dimensional model. \LN
\end{keybox}
\src{\code{H1A-EXACT.md} \S0 and \S4, first item; \S2.3 for the mechanism.}

\subsection{The $N$-scaling of $\beta^{*}$: $\beta^{*}(N)\propto N^{3/2}$}

The $N=3$ measurement leaves one quantitative question: how far does $\beta^{*}$ move with $N$? By a
power, and the power is fixed by the two coefficient laws, both structural.

\emph{The constant term scales like $N$.} In de Wit--'t Hooft form
$J=\prod_{i<j}[\sin(\mu_{ij}/2)/(\mu_{ij}/2)]^{2}$ and $\Phi=2\sum_{i<j}\varphi(\mu_{ij})$ with
$\varphi''(0)=1/12$: there are $N(N-1)/2$ factors, but each carries the weight $(\nu_i-\nu_j)^{2}$, and
\begin{equation}
  \sum_{i<j}(\nu_i-\nu_j)^{2} \;=\; N\sum_i\nu_i^{2} \;=\; N\,\mathrm{Tr}\,v^{2} \;=\; \tfrac{N}{2}
  \qquad\Longrightarrow\qquad
  \Phi''(0)=\tfrac16\sum_{i<j}(\nu_i-\nu_j)^{2}=\tfrac{N}{12},
  \label{eq:phipp}
\end{equation}
i.e.\ the $N^{2}$ factors collapse into a \emph{single} $N$. The identity was checked digit by digit for
$N=2,\dots,6$ and at $N=2$ reproduces $\Phi_\rho/\rho|_{\rho=0}=1/6$ of model 1. \LT
Measured: $W_0(2)=0.2051153$, $W_0(3)=0.3453\pm0.001$, $W_0(4)=0.4908\pm0.04$, i.e.\ $W_0\propto N$;
both $N^{0}$ and $N^{2}$ are excluded by the data. \LE

\emph{The quadratic coefficient scales like $1/N^{2}$.} With $G=-\beta\,\partial_v f$ and
$f=(1/N)\operatorname{Re}\mathrm{Tr}(UB)$ --- the $1/N$ normalization of the Wilson action --- one has
$\partial_v f=O(1/N)$, hence $V_0=\operatorname{Var}(\partial_v f)=O(1/N^{2})$; measured
$V_0(2)=0.0110593$, $V_0(3)\approx9.0\times10^{-4}$, $V_0(4)\approx4.9\times10^{-4}$. \LE

Therefore
\begin{equation}
  \beta^{*}(N)=\sqrt{W_0(N)/V_0(N)}\;\propto\;\sqrt{N\big/(N^{-2})}=N^{3/2},
  \label{eq:n32}
\end{equation}
and the comparison with $\beta_c$ runs as follows.

\begin{table}[h]
\centering\footnotesize
\setlength{\tabcolsep}{6pt}
\renewcommand{\arraystretch}{1.1}
\begin{tabular}{@{}r r r r r@{}}
\toprule
$N$ & $\beta^{*}$ (two-term law) & $\beta^{*}$ (bare $N^{3/2}$ law) & $\beta_c$ & $\beta^{*}/\beta_c$\\
\midrule
$2$ & $\mathbf{4.30661}$ & $4.30661$ & $2.2986$ & $\mathbf{1.87}$\\
$3$ & $\mathbf{19.6}$ & $7.91$ & $5.6924$ & $\mathbf{3.45}$\\
\bottomrule
\end{tabular}
\caption{The zero of the two-term law against the deconfinement $\beta_c$, both in the same Wilson
normalization $\beta=2N/g^{2}$ (so the ratio is meaningful). The bare $N^{3/2}$ law, normalized at
$N=2$, is the more conservative column; both entries still exceed $\beta_c$. Numbers from
\code{BETASTAR-N.md} \S2.}
\label{tab:betastar}
\end{table}

\noindent
\textbf{Both groups therefore sit on the same side.} At $N=2$ the two-term law is the closed form
\eqref{eq:delta-beta}: $\beta^{*}(2)=4.30661$ against $\beta_c(\mathrm{SU}(2))=2.2986$, a ratio of
$1.87$. At $N=3$ the same law with the measured coefficients gives
$\beta^{*}(3)=\sqrt{0.3453/9.0\times10^{-4}}\approx19.6$, against
$\beta_c(\mathrm{SU}(3))=5.6924$, a ratio of $3.45$; even the bare $N^{3/2}$ law, the conservative
column, gives $7.91>5.6924$. The ratio $\beta^{*}/\beta_c$ thus grows from $1.87$ to $3.45$ with $N$:
\textbf{the window in which H1 holds opens faster than the scaling region does.} \LE

\begin{keybox}
\textbf{Third reading, and the one that closes this section.} $\beta^{*}(N)\propto N^{3/2}$ puts
$\beta^{*}(3)\approx19.6\gg\beta_c(\mathrm{SU}(3))=5.6924$ and
$\beta^{*}(2)=4.30661>\beta_c(\mathrm{SU}(2))=2.2986$, with $\beta^{*}/\beta_c=1.87\to3.45$. \textbf{H1
therefore holds on the scaling region of both groups.} The exact failure of \eqref{eq:delta-beta} happens
only \emph{outside} the $\mathrm{SU}(2)$ $\beta_c$, where the model is no longer read as a lattice
approximation of Yang--Mills at all. \LE
\end{keybox}

\noindent
\textbf{Why $19.6$ is the conservative entry, not the sharp one.} For $N\ge3$ the term $E[H\mid0]$ stops
vanishing, so $\delta$ carries a positive linear piece $a\beta$ ($a>0$, $a\propto1/N$) on top of
$W_0-V_0\beta^{2}$; the measured $\delta$ is monotone increasing (\eqref{eq:n3}) and a two-parameter fit
returns $V_0(3)<0$, so on the measured range $0\le\beta\le10$ there is no zero at all and the sharp
statement remains ``$\beta^{*}(3)$ does not exist''. Both readings give the same verdict, and the verdict
is quoted from \eqref{eq:n32} and \tablename~\ref{tab:betastar}, not from an extrapolation. \LE
\src{\code{BETASTAR-N.md} \S1.1--\S1.2 ($\Phi''(0)=N/12$, $W_0\propto N$), \S2 ($V_0\propto1/N^{2}$,
$\beta^{*}(N)\propto N^{3/2}$, the $\beta^{*}/\beta_c$ table), \S3 (the $2N/g^{2}$ normalization, without
which the ratio would be meaningless) and \S4 (the $a\beta$ term; the sentence that $19.6$ is the
conservative value, the third term only pushing the zero further out); summarised in
\code{LOOP-RUNBOOK.md} \S11.48, with the three $\delta$ points of \eqref{eq:n3} in \S11.46. Running code:
\code{research/unconditional/w7/betastar_scaling.py}. \textbf{Caveat, again not to be dropped: single
plaquette/patch Monte Carlo, not the full theory}; the $N$-scaling is a statement at the same $a$ and in
the same $\beta$ normalization, and no $a\to0$ limit is taken.}

\subsection{External validation, and one self-correction kept on the record}

\textbf{External validation.} Five numbers agree digit by digit with the archived corpus (relative
difference $<10^{-6}$): $\Gpp(0)=0.205115$, $\Gpp(2.5)=0.135995$, $\Gpp(6)=-0.193019$,
$\Var(1)/\beta^{2}=0.0110593$, $\beta^{*}=4.30661$. \LE
\src{\code{H1A-EXACT.md} \S2.4, against the \code{[P1]} table and the \code{[A0]} line of
\code{research/v3/species/w7/na_plaquette.log}.}

\noindent
\textbf{Self-correction kept on the record.} The first version of the script returned
$\Gpp(2.5)=+0.5427$; the cause was double counting the second-order term of $\sin(\rho/2)u$, namely
$s\cos(r/2)/(2r)$, which vanishes because $u''=0$ and $\rho'=0$. The test is the archived
$E[H]\equiv0$. \LE
\src{\code{H1A-EXACT.md} \S2.4, last paragraph; errata table row 5 of \S\ref{sec:errata}.}

\subsection{The exact window, and the measured one}
\label{sec:window}

The certificate of \S\ref{sec:scalar-cert} is a statement at fixed $\theta$. Two moment inequalities
extend it to a region of $(\beta,\theta)$, and the region comes in two grades whose difference is one
unproved monotonicity.

\begin{keybox}
\textbf{Theorem (the analytic window).} On the $\xi_{3}=0$ slice of model~1, with
$w(\rho)=\sin(\rho/2)/\rho$ and $\mu\propto e^{-\beta(1-c\cos(\rho/2))}\sin^{2}(\rho/2)\rho^{-1}
\mathrm{d}\rho$,
\[
\Var(G)\le E[H]\iff 2\beta s^{2}\Var_{\mu}(w)\le c\,\langle w\rangle
\qquad\text{and}\qquad
\beta\sin^{2}\tfrac{\theta}{2}\le\cos\tfrac{\theta}{2}\ \Longrightarrow\ \Var(G)\le E[H],
\]
hence $H1$ holds for every $\beta\ge\beta_{0}(\theta)=\sin^{2}(\theta/2)/\cos(\theta/2)$: at
$\beta=2.5$ this covers $\theta\le0.388\pi$ and at $\beta=6$ it covers $\theta\le0.256\pi$. \LT
\end{keybox}
\src{\code{Q04-WEYL.md} \S4.1 (the pointwise ratio $H=\tfrac12\cot(\theta/2)G$ makes the comparison one
inequality, and $0\le w\le\tfrac12$ gives the sufficient condition); \code{R02-CRIT.md} \S3 (the same
implication, now with a proof of the moment inequality rather than an observation about the range of
$w$); \code{R09-FLUCT.md} \S4 (third derivation, same threshold); \code{docs/paper-v2/ch19-h1-status.tex}, subsection ``one window, at two grades'', for the two grades
side by side.}

\noindent
The measurement widens this by a factor $2.3$ in $\theta$ at $\beta=2.5$ and $3.1$ at $\beta=6$, and
the widening is exactly the constant that the stronger of the two moment inequalities buys. With $R:=\langle w\rangle/\Var(w)\ge2$ proved
unconditionally, the numerical minimum over $\gamma=\beta c\in[0,\infty)$ is
\begin{gather}
R_{\min}=\frac{2N_{1}Z}{N_{2}Z-N_{1}^{2}}=32.3177375272,\qquad
Z=\tfrac12\mathrm{Cin}(2\pi),\notag\\
N_{1}=\tfrac34[\mathrm{Cin}(3\pi)-\mathrm{Cin}(\pi)],\qquad
N_{2}=\mathrm{Cin}(4\pi)-\mathrm{Cin}(2\pi),
\label{eq:Rmin}
\end{gather}
attained at $\gamma=0$ (a boundary value, hence exact in $\gamma$), so
\begin{gather}
\beta\le\beta_{c}(\theta):=\tfrac12 R_{\min}\frac{\cos(\theta/2)}{\sin^{2}(\theta/2)}
=16.15887\,\frac{\cos(\theta/2)}{\sin^{2}(\theta/2)}\notag\\
\Longrightarrow\qquad
\theta\le0.904\pi\ \ (\beta=2.5),\qquad \theta\le0.786\pi\ \ (\beta=6).
\label{eq:betac}
\end{gather}
\textbf{That is \LE, not \LC}: the only step missing between \eqref{eq:betac} and a theorem is the
monotonicity of $\Var(w)$ in $\gamma$ (equivalently $R\ge32.32$), and with the proved bound $R\ge2$ the
widening shrinks to $36\%$--$39\%$ ($\theta\le0.527\pi$, $0.356\pi$), which is also measured because the
better constant there is $R_{\min}=4.017$ on the multi-period window.
\src{\code{R05-WCF.md} \S2--\S5 (\eqref{eq:Rmin}: the closed forms at $\beta=0$, the $1/\gamma$
expansion, the $41$-point table, and the minimum at $\gamma=0$), \S6.2--\S6.3 (\eqref{eq:betac} and the
sentence naming the missing monotonicity); \code{R02-CRIT.md} \S4 (the $R_{\min}=4.017$ column and its
dependence on the window $\Lambda$).}

\noindent
Two boundaries travel with the window, and one of them is quantitative enough to close the question.

\textbf{(1) No sharpening of these two legs makes the window $\theta$-uniform} \LN. The criterion is a
product, $\beta\,w(\theta)\ge\tfrac19$ with $w=\cos^{3}(\theta/2)/\sin^{2}(\theta/2)$ and $w(\pi)=0$, so
$\sup_{\theta}\beta_{0}=+\infty$ and no reparametrization can help, the product form being preserved;
quantitatively $\beta_{0}\asymp8/(9(\pi-\theta)^{3})$, which reproduces the tabulated
$2.87\times10^{4}$ at $\theta=0.99\pi$ to three digits. \textbf{(2) At $\theta=\pi$ the premises
collapse} \LM: $c=0$ forces $B=\beta c\equiv0$, the measure is exactly Haar, $\Var=\beta^{2}V_{0}$ is
unbounded, $V_{\infty}=+\infty$ and $c_{\infty}=0$ with $E[H_{vv}]\equiv0$ --- the variance leg is
vacuous and the mean-side leg is vacuous. This is the same slice on which
$\Gpp=0.2051153-0.0110593\beta^{2}$ is exact, so certificate and counterexample are the two faces of
one degeneration. One honest widening on the other side: $\Var\le\beta^{2}V_{0}$ for all $\theta$ gives
$H1$ for $\beta\lesssim3.9$ at every $\theta$ (\LC, since $\inf\Wf$ is a table infimum); what remains
uncovered is the bounded rectangle $\{\theta>0.8335\pi\}\times\{\beta\in(3.9,6]\}$, and it contains
$(\theta=\pi,\beta>4.30661)$, i.e.\ a point where H1 is \emph{false} rather than unproved.
\src{\code{P20-THETA.md} \S4.1 (the product form and the impossibility), \S4.2 (the anchors
$\beta_{0}=2.50$ at $\theta_{\max}=0.7792\pi$ and $6.00$ at $0.8335\pi$), \S3 (the premise-collapse
table), \S4.4--\S4.5 (the small-$\beta$ leg, $\beta_{\rm unif}\approx3.9$, and the residual rectangle,
which contains a false point). \textbf{Note on a numerical coincidence}: \eqref{eq:betac} at $\beta=6$
gives $0.786\pi$, within $0.9\%$ of the $\beta_{0}$-route value $0.7792\pi$; the two routes compute
different functionals ($\Delta_{\mathrm{dyn}}\ge0$ versus $\Gpp\ge0$ with $\Wf$ included), so agreement
is not required and the closeness is evidence about the geometry of the slice, not a derivation
either way.}

\subsection{The lattice ledger: $L=2$ and $L=3$}
\label{sec:lattice-ledger}

The measured side of the same quantity is a Monte Carlo on the four-dimensional model, with the split
observation $S_{W}$ versus $\Phi$ maintained per sample so that
$\Delta_{\mathrm{dyn}}+\Wg+X-\Gpp$ is checked on every configuration (sum check $\le1.6\times10^{-6}$
at $\beta=6$).

\begin{table}[h]
\centering\footnotesize
\setlength{\tabcolsep}{6pt}
\renewcommand{\arraystretch}{1.1}
\begin{tabular}{@{}llll@{}}
\toprule
volume & $\beta$ & $\Gpp$ & statistics\\
\midrule
$L=2$ & $1$   & $+0.2835\pm0.0049$ & $27.0\sigma$ (jk)\\
$L=2$ & $2.5$ & $+0.4116\pm0.0421$ & $4.0\sigma$ (jk); $+0.3942\pm0.0258$ at $6.0\sigma$ on a
$1.5\times10^{4}$-sweep run\\
$L=2$ & $4$   & $+0.5108\pm0.0939$ & $2.3\sigma$ (jk)\\
$L=2$ & $6$   & $+0.6018\pm0.0589$ & $10.2\sigma$; five estimators span $4.1$--$10.2\sigma$\\
$L=3$ & $2.5$ & $+0.3396\pm0.0585$ & $5.80\sigma$ (jk) / $2.24\sigma$ ($\tau$)\\
\bottomrule
\end{tabular}
\caption{$\Gpp=E[\kappa]-\Var(F)$ on the lattice, mode $A$, $\|v\|=1$, finite-difference derivatives at
$h=10^{-3}$. Every entry is \LE: one volume window, one mode, no $a\to0$.}
\label{tab:ledger}
\end{table}

\noindent
Three readings. \emph{(i) The sign is uniform in $\beta$ on this window}: the four $L=2$ central values
increase monotonically $0.2835\to0.4116\to0.5108\to0.5913$ with no interior minimum and no crossing, and
$16/16$ chains are positive (weakest $+0.1338$ at $\beta=4$). The mechanism is a competition that does
not flip: $D/S$ rises $0.505\to0.787\to0.848\to0.889$, approaching $1$ without crossing, so
$\Gpp=E[\kappa](1-D/S)>0$ with a shrinking margin. \textbf{The extrapolation beyond $\beta=6$ is an
extrapolation}: a quadratic fit puts the zero near $17.3$, a cubic near $14.9$, and neither is a result.
\emph{(ii) At $\beta=6$ the statistics suffice only after the estimator is audited}: $4$ chains
$\times\,5\times10^{3}$ sweeps gave $+0.5913$ at $1.3\sigma$ because $\tau_{K}=100$--$240$; $8$ chains
$\times\,2\times10^{4}$ sweeps give $+0.6018$ with jk25 $10.2\sigma$, eight-chain scatter
$9.2\sigma$, long-block batch means $8.2\sigma$, $\tau_{z}$-corrected $7.3\sigma$ and the old
$\tau_{F}$ formula $4.1\sigma$; $8/8$ chains positive, sign test $p=3.9\times10^{-3}$,
$\chi^{2}/\mathrm{dof}=1.22$. The single estimator below $5\sigma$ is a formula mismatch: it applies
$\tau_{F}$ to the fluctuation of the sample variance, whose correlation time is $\tau_{z}=8.7$--$14.1$,
and treats a fourth moment as Gaussian. \emph{(iii) There is no volume dependence}: at $\beta=2.5$,
$L=2$ gives $+0.3942$ and $L=3$ gives $+0.3396$, a difference of $-0.055\pm0.064$ ($0.85\sigma$), while
$\Var(F)$ agrees to $0.7\%$ ($1.535$ against $1.5458$). The $L=3$ number also settles a sampler bug: the
earlier $-0.177\pm0.062$ came from a chain with no fundamental-domain constraint, which wandered onto
the repeated flat chart where $\Phi\sim2\log(r/2)$ supplies no restoring force; with the hard rejection
$|x_{l}|<2\pi$ the variance falls from $1.8917$ to $1.5458$ and the sign flips.
\src{\code{P19-BETA-AT-L2.md} \S0 and \S2; \code{Q13-BETA6.md} \S0--\S3; \code{P01-L3.md} \S0--\S1;
\code{O24-L2-PUSH.md} \S2 and \S5; \code{LOOP-RUNBOOK.md} \S11.160--\S11.161 for the combined
statement.}

\subsection{Where exactness stops: $143$ dimensions and seven}
\label{sec:exactness-stops}

\textbf{The volume.} At $L=2$ there is no finite-dimensional reduction. The dimension count is exact:
$192=4L^{4}\cdot3$ link components, minus $48$ for the gauge group, minus one for the conditioning
$\langle v,X\rangle=a_{0}$, leaves
\begin{equation}
192-48-1=143 ,
\label{eq:143}
\end{equation}
and $143\gg20$ is the criterion for ``no interval arithmetic''. Symmetry is a theorem, not a hope:
$W_{4}$ (order $384$) and the space group (order $6144$) are finite, so every orbit is a finite set and
$\dim(M/G)=\dim M$; what a finite group buys is the measure factor $1/6144$ and a three-dimensional
invariant tangent space at the flat point, which is the free field where $\Gpp\equiv0$ and the acid test
fails. The per-link one-dimensional factorization that made the certificate possible does not exist
here, because $S_{W}$ is a sum over plaquettes; the by-parts pairing giving $NS=1/4$ is an accident of
one plaquette. What is available instead is a statistical certificate whose price is explicit: a
bounded-part-plus-tail-volume Hoeffding bound needs $N\approx3\times10^{8}$ samples ($\approx35$ days
on one core). One exact simplification does exist, at the corner $a_{0}=16\pi$:
$\rho_{v}\propto\varepsilon^{191}$ and $\Gpp=+191/\varepsilon^{2}\to+\infty$ --- $Q08$'s
$+3/(2\pi-a)^{2}$ promoted to $L=2$ --- but the corner is $16\pi$ from the conditioning point. \LN
\src{\code{R04-L2CERT.md} \S0--\S2 (the count, the rank checks, the finite-group theorem, the orbit
counts, and the $12$-dimensional toron decomposition $192=45+12+135$), \S3 (the cost of a rigorous
statistical certificate), \S4.1 (that $\theta=\pi$ does not simplify at $L=2$, and the $3n-1$ corner law
with $n=64$); \code{Q08-LARGEA.md} \S2.1 for the single-plaquette version.}

\noindent
\textbf{The group.} $\mathrm{SU}(3)$ is measured, not proved, and the two numbers must be quoted
together because their signs differ:
\begin{equation}
\mathrm{gap}\big|_{\text{weight }|\Delta|^{2}e^{-S_{W}}}=+0.10172408,\qquad
\mathrm{gap}\big|_{\text{pure Haar}}=-0.23580249
\qquad(\theta=\pi,\ \beta=2.5,\ a=0).
\label{eq:su3gap}
\end{equation}
Two earlier attempts are dead, for reasons worth keeping. The first inferred $S\equiv D$ from
$S_{W}\equiv\beta$ on the slice; that is invalid, because $S_{W}\equiv\beta$ is a \emph{value} and
$G=\partial_{v}S_{W}$ is a derivative --- verified: $S_{W}-\beta=0.0$ at three points while
$G=\beta sg(u)=1.22167/0.97087/0.24784$, matching finite differences to eighteen digits. The second
error is the choice of $\Phi$: it used $-\log|\Delta|^{2}$ where the corpus object is
$\Phi=\sum_{i<j}g(d_{ij})$, the two differing by $\psi=\sum_{i<j}\log d_{ij}^{2}$; with the wrong choice
the identity $E[\Phi'^{2}]=E[\Phi'']$ holds exactly ($48=48$ in exact rationals, width
$1.2\times10^{-119}$) and forces $S-D=-\Var(G)-2\mathrm{Cov}\le0$. The obstacle to a certificate is
dimension, not precision: the $a=0$ slice is seven-dimensional (one eigenvalue direction times a
five-dimensional flag), two-dimensional Weyl integration is exact only for class functions while $G$
and $\Phi'$ are not, and $\Phi$ contains $\log$ and $\cot$/$\csc^{2}$ terms that destroy the
trigonometric-polynomial route that worked at $N=2$. \LE{}/~\LM
\src{\code{R12-SU3CORR.md} \S0--\S3 (both numbers, the $G\neq0$ verification, the two errors, the exact
$48=48$, and the seven-dimensional obstacle), \S4 (the verdict); \code{R11-SU3ADJ.md} (the adjudication
that the $S\equiv D$ inference is invalid); \code{R08-SU3CERT.md} (the file under adjudication, kept on
the record).}

\subsection{What this section does not claim}

\begin{itemize}
\item It does \textbf{not} claim that the paper's H1 (fixed physical $\beta$, $(L,a)$-uniform, every
      coarse mode) has been refuted. What was refuted is the stronger reading ``H1 for all $\beta$ /
      all vacua''. \LT
      \src{\code{H1A-EXACT.md} \S0 and \S4, first item.}
\item It does \textbf{not} claim that $\delta>0$ yields an atom. In 2D YM and in the P4 kernel the
      positive $\delta$ is carried entirely by the Haar term, while at $\beta=0$ one has
      $\Gpp=0.2688$ with $f_1^{2}=0$: a positive $\delta$ is a ceiling, not an atom. \LT
      \src{\code{H1A-EXACT.md} \S2.5, last paragraph, quoting \code{ch15-uvface.tex}:236--241; the
      difference between $\Gpp=1/(2\,\mathrm{Cin}\,2\pi)=0.20511530$ and $0.2688$ is the subject of
      \code{DS-RATIO.md} (errata row 17).}
\item It does \textbf{not} claim that the $\delta$'s above have been scaled to $a\to0$: all of them
      are surfaced bounds in lattice units. \LM
      \src{\code{H1A-EXACT.md} \S4, final paragraph (``not closed: an honest statement'').}
\end{itemize}

% =====================================================================
%  §4  Five conditions $\to$ one target
% =====================================================================
\section{Five conditions $\to$ one target: only $C5$'s rate uniformisation survives}
\label{sec:five}

The paper lists five requirements the ``machine'' must meet (non-perturbative, sign-controlling lower
bound, non-circular, acid-tested, $(L,a)$-uniform). This section is pure bookkeeping: the five are
\emph{not} five independent axes, and scoring them as five manufactures free progress. After the
individual verdicts exactly one survives, and it has a name.

\subsection{The five, one by one}

\begin{table}[h]
\centering\footnotesize
\setlength{\tabcolsep}{4pt}
\renewcommand{\arraystretch}{1.12}
\begin{tabularx}{\textwidth}{@{}l l X l@{}}
\toprule
\textbf{Condition} & \textbf{Grade} & \textbf{Family of the immune graph / ceiling} & \textbf{Source}\\
\midrule
$C1$ non-perturbative & \LN{} (class level) & families 46/47/48/53/54/55; ceiling \textbf{R6}: Balaban stops at
$g^{2}=1=\tH/(2N)$, short by $2.47\times$ (physical point) / $18\times$ (lower edge) & \code{ch13}:105--119\\
$C2$ sign-controlling bound & \LN{}+\LM & families 49/50/56/57, plus \code{OS-positivity} $\not\Rightarrow$
exponential decay; no row in the ceiling table, the ceiling $=$ H1 itself & \code{ch13}:494--508\\
$C3$ non-circular & \LT{}+\LM & family 47 (presupposing the conclusion is one of the death modes); form
satisfied, content (isolation $s_2$) not computable & \code{ch13}:509--513\\
$C4$ acid test & \LE & family 46-acid branch/49/51; the instrument idles (it quantifies over the empty set) & \code{ch13}:514--536\\
$C5$ $(L,a)$-uniform & \LM & the whole R1--R6 table (built for it); best row \textbf{R4} (TRG, HS norm) $2.9\times$ & \code{LOOP-RUNBOOK}:363--372\\
\bottomrule
\end{tabularx}
\caption{The five conditions as they stand. Four already have family-level verdicts; the fifth has a
verdict on its import face only, and none on its residual face.}
\label{tab:five}
\end{table}

\src{The conditions are quoted verbatim from \code{docs/paper-v2/ch13-machine.tex}:12--30, which is
word for word the list of \code{docs/PROOF-SKELETON.md}:74--76 (so there is no conflict of standards);
family numbers from \code{research/v3/species/H1B-FIVE.md} \S1--\S2 and
\code{LOOP-RUNBOOK.md}:1118--1135, :1546--1551.}

\noindent
One erratum that must be recorded: the $8.9\times$ of $R1$ has been corrected to $19$--$20\times$
(the paper itself writes ``an earlier figure of $8.9\times$ made the same conversion slip''). \LT
\src{\code{docs/paper-v2/ch13-machine.tex}:194--195; \code{H1B-FIVE.md} \S C5; the root cause of the
conversion is \code{BETA-CONVENTION.md} and errata rows 4 and 10 of \S\ref{sec:errata}.}

\subsection{Logical relations: one implication, one restatement, two filters}

\begin{enumerate}[label=(\arabic*)]
\item \textbf{$C5\Rightarrow C2$.} $C2$ is ``produce a lower bound
      $m_{\mathrm{eff}}^{2}(k\to0)\ge\lambda>0$''; $C5$ is ``that bound is uniform in $(L,a)$ and
      survives the continuum limit''. $C5$ is $C2$ \emph{plus axis qualifications}: every tool that
      satisfies $C5$ satisfies $C2$, and not conversely. \LC
      \src{\code{ch13}:17--20 and :27--29; \code{H1B-FIVE.md} \S4.1.}
\item \textbf{$C2$'s window coordinate is a restatement.} The paper says the ratio is ``not an
      independent half'' and that its upper half ``is a floor statement inheriting H1''. Hence the
      window supplies no independent half of $C2$. \LT
      \src{\code{ch13}:494, :504; \code{H1B-FIVE.md} \S4.2.}
\item \textbf{$C3$ and $C4$ are filters, not propositions.} The originals read ``must not assume''
      and ``must \emph{fail}'': neither can be refuted nor supported, only violated by a candidate. \LT
      \src{\code{ch13}:21--26; \code{H1B-FIVE.md} \S4.3.}
\item \textbf{$C1$ follows from $C5$ inside the setting (not a purely formal implication).} At
      $\beta\approx2.5$, $g\approx1.5$ there is no small parameter, so any bound uniform in $(L,a)$
      there cannot be the product of a convergent expansion. \LC
      \src{\code{ch13}:14--16, :105--119; \code{H1B-FIVE.md} \S4.1, last paragraph.}
\end{enumerate}

\begin{keybox}
\textbf{An accounting error, established.} The corpus scores the five as independent axes
(\code{BC-BETHE.md}:193, the five-condition scorecard in \S5; \code{LF-LASTFOUR.md}:79/130/184/273,
``1.5/5'', ``1/5'', ``2/5''; \code{MG-META.md}:240/254/266; \code{SDP-CERT.md}:295), while the same
corpus also writes that the five are several spellings of one object. Treating $C1$/$C3$ as
independent scoring axes manufactures free progress: any formal construction automatically passes
$C1$ (it does not expand) and $C3$ (it presupposes nothing), so ``$2/5$'' may equal zero content. \LN
\end{keybox}
\src{\code{H1B-FIVE.md} \S4.4; \code{LOOP-RUNBOOK.md} \S11.33 (lines 1739--1744).}

\noindent
\textbf{The correct convention.} There is one substantive axis (the single lower bound that $C5$
contains together with $C2$); $C3$/$C4$ are scored \textsc{pass}/\textsc{fail}; $C1$ follows from
$C5$ inside the setting and is listed separately as an \emph{admission} condition. \LT
\src{\code{H1B-FIVE.md} \S4.4, last paragraph.}

\subsection{The only target: one name, five spellings}

\begin{keybox}
\textbf{The only substantive target} $=$ the \textbf{rate uniformisation} of $C5$
$=$ the corpus object $d5$\,/\,$(CP)$
$=$ ``the mechanism that transports a rate in lattice units into a physical rate'' (the
\textbf{cross-scale transport machine} of \code{L3})
$=$ $p=1$ as decided in \code{A7} (the exponent must equal $1/(2b_0)$)
$=$ the contradiction of \code{t3}: ``$ma\ge r$'' and ``$\xi/a\to\infty$'' point in strictly opposite
directions. \LM
\end{keybox}
\src{\code{H1B-FIVE.md} \S3; \code{LOOP-RUNBOOK.md} \S11.34 (lines 1746--1759): five independent
routes (\code{t3}/\code{l3}/\code{a7}/\code{h1b}/\code{BC-BETHE}) converge on the same name. The paper
itself marks it \textbf{untouched} \code{[M]} in \code{ch13}:171--173.}

\noindent
Why it is not ``one more constant to tune'': the chain is $\mathrm{RP}\Rightarrow \mathrm{KP}(a)
\Rightarrow$ exponential clustering $\Rightarrow$ KL, and what it delivers is ``gap $\ge a$'' where $a$
is in \emph{lattice} units, while near the root $\xi/a\approx0.5$--$1.3$ shows \textbf{no scale
separation}. Turning ``gap $\ge a$'' into ``gap $\ge c\Lambda$'' is not a matter of a better constant
but of a mechanism that changes units. \LM
\src{\code{ch13}:170--173; \code{H1B-FIVE.md} \S3, first item.}

\subsection{What it lacks (item by item, no vague language)}

\begin{enumerate}[label=(\arabic*)]
\item \textbf{A mechanism transporting a lattice rate into a physical one}: $\xi/a\approx0.5$--$1.3$,
      no scale separation. \LM
\item \textbf{A replacement for that single triangle inequality}: in $\rho=N/2-C^{*}N\beta$ the $N/2$
      is the \emph{exact} $\mathrm{SU}(N)$ Einstein constant (zero slack), while the remaining
      $8(d-1)$ sit entirely inside one triangle inequality whose configuration supremum $C^{*}$ is
      saturated by the maximal centre vortex ($d=2$: $8.000000$; $d=4$: $[16,24]$) $\Rightarrow$
      replacing it \emph{is} H1. \LT
      \src{\code{LOOP-RUNBOOK.md}:368, :377; \code{ch13}:186--197.}
\item \textbf{The continuum limit}: $1/32$ is a number obtained after converting $\beta_A$ into
      $\beta_W$; it is not a family that can be pushed to $a\to0$. \LM
      \src{\code{ch13}:192--195; \code{app:erratum-radius} records the two old mistakes $75\times$ and
      $1/24$.}
\item \textbf{Measurability (not computing power)}: $L=4$ gives $\rho^{2}=41.9\pm46.0$ (error larger
      than the value) and $L=6$ shows no plateau; the free-field control
      $\rho^{2}_{\mathrm{free}}=2<4$ has \emph{no} mass gap, so the criterion cannot separate free from
      interacting. \LE
      \src{\code{ch13}:646--658; \code{GEVP-4D.md}.}
\item \textbf{The time-integrability half}: after the BD books are fixed, the gap moves to an
      $L$-uniform bound on $\|\operatorname{Cov}_{\mu_{t,x}}(\nabla S_W)\|_{\mathrm{op}}$ at
      intermediate $t$, plus integrability in time; both endpoints are proved ($t=0$ is exactly $0$;
      $t=\infty$ is strictly $\le6\beta$). This is the \emph{narrowest} gap of the campaign: it
      contains no volume. \LM
      \src{\code{LOOP-RUNBOOK.md} \S9.4, last paragraph (lines 530--535).}
\end{enumerate}

\subsection{The cheap side is already closed: a quantitative no-go}

To obtain $C5$ from the arithmetic / polymer side one needs the per-plaquette weight decay rate to
exceed the surface entropy. The per-plaquette rate of $w_\beta(S)$ is
$c(\beta)=-\log u=\kappa+\sigma a^{2}/\pi$, and the entropy exponent of plaquette clusters is
$\kappa\in[1.4012,\,4.3307]$ (the lower bound $\kappa_{\min}=1.40$ of $q4$ is an effective bound). At
the physical coupling
\begin{equation}
  c(2.5)=0.679,\qquad \frac{c(2.5)}{\kappa_{\min}}=\frac{0.679}{1.40}\approx0.49 ,
\end{equation}
i.e.\ \textbf{the decay rate is only half the entropy} $\Rightarrow$ the polymer expansion diverges at
the physical coupling. \LN
\src{\code{Q4-LOGCONVEX.md}, as reported in \code{LOOP-RUNBOOK.md} \S11.38 (lines 1825--1849); the
interval for $\kappa$ is from \code{Q5-KAPPA.md} and \S11.40.}

\noindent
Moreover the ``surviving branch'' of that criterion is circular: $c>\kappa \iff \sigma>0 \iff$ area
law. \LN
\src{\code{LOOP-RUNBOOK.md} \S11.38(1) (``tautology: yes'').}

\subsection{Twelve roads, each with its number}
\label{sec:twelve}

The negative asset of the same wave is a list of routes that cannot work, each ended by a number
rather than by an impression. Ten were already in force; two are new, and both are of the kind that
matters --- a route that turns out to be the criterion renamed, and a route whose residual turns out
not to be noise.

\footnotesize
\begin{longtable}{@{}>{\raggedright\arraybackslash}p{0.4cm}
  >{\raggedright\arraybackslash}p{3.0cm}
  >{\raggedright\arraybackslash}p{8.4cm}
  >{\raggedright\arraybackslash}p{2.6cm}@{}}
\caption{The twelve roads, each with the number that ends it. All twelve are class-level \LN: in every
row but two the hypothesis is \emph{false}, and in the other two it is satisfied and the channel is
still too narrow.}\label{tab:twelve}\\
\toprule
\# & Road & Why it dies, with the number & Source\\
\midrule
\endfirsthead
\toprule
\# & Road & Why it dies, with the number & Source\\
\midrule
\endhead
\bottomrule
\endlastfoot
1 & Brascamp--Lieb, all directions & the required $\lambda$ is $38$--$44$ times $N/12$ and
$\mathrm{Hess}(S_{W}+\Phi)$ is indefinite at $\beta=2.5$: the hypothesis fails, not the constant &
\code{Q02-BL-PHI.md}\\
2 & ODE / Riccati for $\Gpp$ & $\Gpp$ does not close: $D'\le S'$ is \emph{equivalent} to $G'\ge0$
($\max|\mathrm{gap}+G'|=2\times10^{-16}$), a tautology that cannot be upgraded; and $D/S$ crosses $1$ at
$\beta_{c}\approx8$--$11$, with $\beta=6$ already at $0.941$ & \code{Q06-RICCATI.md}\\
3 & Curvature-dimension, RC-BL & $\mathrm{ess\,inf}\,K_{+}=0$ (the bad set carries $K_{+}\equiv0$;
$K\ge N/3-|\lambda_{\min}|=-0.893$ for $\mathrm{SU}(3)$), and $\sum\theta_{j}$ cannot be made
consistent & \code{Q14-RCBL.md}\\
4 & Sum of squares & the direction-by-direction pairing needs
$\Var\le N/12-\langle|\lambda_{-}|\rangle=-0.0367<0$ with $\langle\lambda_{-}\rangle=-0.2033$: false for
every $\Var\ge0$, pure arithmetic & \code{Q12-SOS.md}\\
5 & Cauchy--Schwarz & loose by $2.6\times$ on the upper side and $40$--$125\times$ on the lower side &
\code{Q09-COVLOW.md}\\
6 & Spectral conditions & $\lambda_{k}$ cancels term by term in $E[H]-\Var(G)$; the named criterion
$\Lambda\preceq I$ holds on $0.00\%$ of configurations ($\lambda_{\max}=\beta=2.5$) while the target
holds and is an equality & \code{Q15-SPECTRAL.md}\\
7 & Two-halves decomposition & the $\Phi$-only margin is \emph{exactly} $0$ unconditionally ($1/3$
against $1/3$, by parts at machine precision), and on the interacting side $\Phi$ turns negative at
$a>1.588$ & \code{Q01-PHI-ONLY.md}\\
8 & The $\beta_{0}$ window's small-$B$ bridge & the claimed bound
$\Var\le\beta^{2}s^{2}(1/2-1/\pi)^{2}/4$ is broken by a counterexample inside its own window: at
$B=0.5$, $\Psi=0.082919>0.2971B^{2}=0.0743$ & \code{P20-THETA.md} \S6.5\\
9 & Pure-$\Phi$ quantities as the witness & they do not vanish on the free field:
$\Var(\Phi')=2.341\times10^{-3}$, i.e.\ $65\%$ of the interacting value, and $E[\Phi'']-N/12$ is
identical at $\beta=2.5$ and $10^{-6}$ & \code{P21-REBUILD.md}\\
10 & Statements uniform in $a$ & the negative region grows monotonically in $\beta$: width
$1.09\to2.26\to3.57\to4.36\to5.22\to5.66$, depth $-0.016\to-22.7$ for $\beta=2\to40$ &
\code{Q08-LARGEA.md}\\
11 & Curvature fluctuation $\beta\Var(\partial_{t}h)\le E[\partial_{t}^{2}h]$ & it is the criterion
renamed ($S_{W}=\beta h$ makes both sides the same ratio) and it fails at $\theta=\pi$ for every
$\beta>0$: LHS $=0.027648>0=$ RHS, because $H_{vv}\equiv0$ is a pointwise identity &
\code{R09-FLUCT.md}\\
12 & Affine reduction $H=K(a)G+C(a)$ & not an identity: $W(G,H)=O(a)\not\equiv0$,
$W(G',H')\not\equiv0$, the two channels differ $30$-fold ($K_{c}=10.356$ against $K_{s}=0.308$), and the
residual is $O(a)$ with $\mathrm{relres}/a\to1.43\times10^{-3}$, stable under a fourfold quadrature
increase & \code{R07-AFFINE.md}\\
\end{longtable}
\normalsize

\noindent
An accounting point about this table, because the temptation with such a list is to count it. \emph{The
twelve deaths are not twelve independent pieces of evidence.} Rows 1--6 have a common root, stated in
the corpus as $E[H_{vv}]|_{a\neq0}<0$: every one of the six is a tool that can produce
``variance $\le$ curvature'' only if that mean curvature can be made positive, and at $a\neq0$ it
cannot. Rows 7--12 are structural variants of the same obstruction rather than new obstructions. What
the table therefore establishes is one negative statement with twelve witnesses: \textbf{the missing
object is a directional, conditioned, sign-carrying lower bound on $E[H_{vv}]$ at $a\neq0$, and no
classical variance--curvature tool delivers one.} \LN
\src{Rows 1--10 are \code{LOOP-RUNBOOK.md} \S11.158 (the ten-row table with the same numbers) with
\S11.136 and \S11.150 (BL and the common root), \S11.149 (SOS) and \S11.152 (the corrected constant
$C=3.95$, which supersedes both $1.7454$ and $10.05$); rows 11--12 are that day's additions,
\code{R09-FLUCT.md} \S2--\S4 and \code{R07-AFFINE.md} \S3--\S6; row 8's counterexample is
\code{P20-THETA.md} \S6 item 5, with the same arithmetic caught independently in
\code{Q16-CONFLICT.md};
\code{docs/paper-v2/ch19-h1-status.tex}, subsection ``twelve roads that end'', carries the same table
inside the ledger.}

\subsection{Conclusion of this section}

\begin{itemize}
\item Five conditions $\to$ one substantive axis $\to$ one named object: the rate uniformisation of
      $C5$, i.e.\ $d5/(CP)=p=1$. \LM
\item That object is \textbf{not} a shape statement: it asks first for ``gap $\ge a$'' (already
      delivered on the strong-coupling side) and then for a replacement of $a$. It therefore does not
      conflict with the \LN{} of \S\ref{sec:position} --- the latter closes ``deducing the gap from the
      shape'', the former wants ``transporting an existing gap''. \LT
\item It remains untouched by anyone: the paper says \code{untouched} \code{[M]}, and every closure
      of the day happened elsewhere. \LM
\end{itemize}


% =====================================================================
%  §5  Closure of the string-tension side
% =====================================================================
\section{Closure of the string-tension side: $\sigma>0$ does not entail $m>0$}
\label{sec:tension}

This section closes what looks like the cheapest route: the area law is already proved at strong
coupling, $\sigma>0$ is the standard confinement indicator, and ``$\sigma$ and $m^{2}$ have the same
dimension'' seems to leave only a dimensional analysis. Four steps later the route is closed:
$\sigma$ on the lattice is a power-divergent $1/a^{2}$, $\sigma=2g^{2}Nm^{2}$ can only be scheme
bookkeeping, the proved implication runs the other way, and the sector that produces the area law
carries \emph{its own} gapless direction.

\subsection{Step 1: the area law is indeed a solution of the two-variable loop equation}

Insert $W=e^{-\sigma A}$ into the Makeenko--Migdal equation
\begin{equation}
  \partial_\mu\frac{\delta W}{\delta\sigma_{\mu\nu}(x)}
  \;=\; g^{2}N\oint_{C}\dd y_\nu\,\delta^{4}(x-y)\,W(C_{xy})W(C_{yx}),
  \label{eq:MM}
\end{equation}
and on the planar circle family $\theta$ drops out exactly: the two circular segments satisfy
$A_1+A_2=\tfrac12R^{2}\cdot2\pi=\pi R^{2}$ as an identity, so $G\equiv W$. The residual
\begin{gather}
  \rho(R,\theta):=\frac{\text{LHS}}{\text{RHS}}=\frac{\sigma}{2g^{2}N}
  \ \equiv\ 0.500000000000 ,\notag\\
  \rho(5)/\rho(1)=1.000000000000 ,\notag
\end{gather}
has $R$-power $0$ and $\theta$-mismatch $0$ (an exact identity, not a small number); the
one-dimensional reduction on the same $(R,\theta)$ grid runs $5.4366\to8.821\times10^{14}$ (ratio
$1.622\times10^{14}$). Hence \LT
\begin{equation}
  \sigma \;=\; 2g^{2}N\,\mu^{2},
  \qquad \mu = \text{contact-term / coincident-point scale}.
  \label{eq:sigma-mu}
\end{equation}
\src{\code{RV-2-AREALAW.md} \S1--\S5; witness \code{w7/rv2_arealaw.py} / \code{.log} (PART 1 three
geometry checks, PART 3 twenty residuals, PART 5 the dimensional table).}

\noindent
\textbf{Its limitation must be recorded at the same time}: in the planar family $A_1+A_2\equiv\pi R^{2}$
turns the equation into a pointwise \emph{algebraic} relation and forces $G_\theta\equiv0$.
\textbf{The planar family does not test the functional form of the area law, only its coefficient.} \LT
\src{\code{RV-2-AREALAW.md} \S6, the ``structural theorem'' item.}

\subsection{Step 2: on the lattice $\sigma\propto1/a^{2}$, exactly}

Putting \eqref{eq:sigma-mu} back on the lattice needs only arithmetic. With
$\delta^{4}(x-y)\to a^{-4}\delta_{x,y}$ (a dimensionful density), a line integration carrying the link
length $a$, and a one-dimensional curve delta $\delta_C=a^{-1}$: \LT
\begin{equation}
  \delta^{3}_{\perp}(0)=a\cdot a^{-4}=a^{-3},
  \qquad
  \mu^{2}=\frac{\delta^{3}_{\perp}(0)}{\delta_C}=a^{-2},
  \qquad
  \sigma=\frac{2g^{2}N}{a^{2}}\ \propto\ \frac{1}{a^{2}} .
\end{equation}
\src{\code{MU-4-LATTICE.md} \S1--\S3; scripts \code{w7/mu4_lattice.py}, \code{w7/mu4_lattice.log}.}

\begin{table}[h]
\centering\footnotesize
\setlength{\tabcolsep}{7pt}
\renewcommand{\arraystretch}{1.1}
\begin{tabular}{@{}l l l l@{}}
\toprule
\textbf{Normalization} & \textbf{Contact term} & \textbf{$a$-power} & \textbf{Fit ($L=4,6,8,16,32$)}\\
\midrule
$K_A$: dimensionless Kronecker, no $a^{4}$ & $1$ & $0$ & $p=0.0000$\\
$K_B$: $\delta^{4}\to a^{-4}\delta_{x,y}$ with line measure $a$ & $a^{-3}$ & $-3$ & $p=+3.0000$\\
$K_C=K_B/\delta_C$: divided by the curve delta $a^{-1}$ & $a^{-2}$ & $-2$ & $p=+2.0000$\\
\bottomrule
\end{tabular}
\caption{The three normalizations separated by varying $a$. One has
$\sigma_C a^{2}=\sigma_B a^{3}=2g^{2}N$ identically. At $a=1$ all three agree numerically --- which is
why the original check had zero resolving power.}
\label{tab:apower}
\end{table}

\noindent
Two independent numerical criteria point the same way (fixed physical volume $La=1$): \LE
\begin{itemize}
\item $R:=\sigma a^{2}/e^{-1/(2b_0g^{2})}=156.2/205.4/254.0$ ($L=4,6,8$), fitted by
      $R\propto L^{+0.700}$ --- \emph{running}, not constant;
\item $\sigma/\Lambda^{2}=351/629/982$ (diverging like $L^{2}$) --- $1/\Lambda^{2}$ is excluded.
\end{itemize}
\src{\code{MU-4-LATTICE.md} \S3.3 (with the row-by-row table of $x=1/(2b_0g^{2})$, $g^{2}$,
$\beta=2N/g^{2}$).}

\begin{keybox}
\textbf{A methodological accident that must be recorded.} Before the lattice arithmetic was done, a
broadcast vote decided \textbf{5:1} that $\mu^{2}$ should be $1/\Lambda^{2}$. The exact lattice
arithmetic overturned it: $\varepsilon=a$ rather than $1/\Lambda$, so $\sigma\propto1/a^{2}$ and the
continuum content is zero. A vote is not a criterion. \LN
\end{keybox}
\src{\code{LOOP-RUNBOOK.md} \S11.25--\S11.26 (the 5:1) against \S11.30 (the overturning);
\code{MU-4-LATTICE.md} \S0; the same family is recorded in \S11.24 (``the $\varepsilon=1/\Lambda$
conjecture is refuted'') and \S11.20.}

\subsection{Step 3: $\mu^{2}$ is a renormalization constant, not a mass}

The $\mu^{2}$ of \eqref{eq:sigma-mu} occupies the dimensional slot of $(\text{mass})^{2}$, but on the
lattice it is forced to $1/\varepsilon^{2}$, and $\varepsilon$ is a \textbf{prescription}: \LT
\begin{center}\footnotesize
\begin{tabular}{@{}l l@{}}
\toprule
Prescription & $\mu^{2}$\\
\midrule
lattice ($\varepsilon=a$) & $1/a^{2}$\\
dimensional regularization & $0$ (the power divergence is removed wholesale, no pole)\\
point splitting & $(4\pi)^{-3/2}/\varepsilon^{2}$\\
heat kernel & $\Lambda^{2}/(4\pi)$\\
\bottomrule
\end{tabular}
\end{center}
Four prescriptions, four values $\Rightarrow$ \textbf{prescription dependence holds} $\Rightarrow$
$\mu^{2}$ is a renormalization constant. \LT
\src{\code{MU-1-PRESCR.md} \S1, \S7(3), collected in \code{MU-3-RENORM.md} \S0.5 and \S6(0); this
document does not repeat the table, only its conclusion.}

\noindent
The criterion is purely exponential, with no number to hide behind: fixed physical $\sigma$ requires
$\dd\ln(\sigma a^{2})/\dd\ln a=2$ (exactly), whereas the candidate $\mu^{2}\simeq1/a^{2}$ gives
$\dd\ln(\sigma a^{2})/\dd\ln a=1/t\to0$; the exponents $2$ and $1/t$ are of \emph{different type}, so
no common interval can exist. \LT
\src{\code{MU-3-RENORM.md} \S5.2.}

\subsection{Step 4: $\sigma=2g^{2}Nm^{2}$ is scheme bookkeeping only}

Renaming the $\mu^{2}$ of \eqref{eq:sigma-mu} as $m^{2}$ leaves two factors: \LT
\begin{center}\footnotesize
\begin{tabularx}{\textwidth}{@{}l l X@{}}
\toprule
Factor & Origin & Nature\\
\midrule
$2$ & normalization of $\tfrac12\varepsilon_{\mu\nu}$ (i.e.\ $1/(\tfrac12)$) & pure convention, no physics\\
$g^{2}N=\lambda(\mu)$ & the right-hand coefficient of \eqref{eq:MM}; in 4D it runs and is scheme dependent & \textbf{the physical content is here}\\
\bottomrule
\end{tabularx}
\end{center}
The physical ratio $\sigma/m^{2}$ must be RG-invariant (both are spectral quantities), while
$2\lambda(\mu)$ runs; hence, if $m$ is a physical mass, this relation \textbf{can only be a
scheme-dependent bookkeeping identity, never an identity}. Conversely, it is precisely that running
factor that demotes $\mu^{2}$ from ``mass'' to ``renormalization constant''. \LT
\src{\code{MU-3-RENORM.md} \S7; the same section gives the criterion: if $\sigma/m^{2}$ is ever
measured to be a universal $O(1)$ number, then the $\mu^{2}$ here must be renamed away from $m$.}

\begin{keybox}
\textbf{$\sigma>0\not\Rightarrow m>0$.} The literature contains no proved
$\sigma>0\Rightarrow m>0$; and the proved arrow runs the other way: \emph{exponential decay of
correlations $+$ nontrivial centre $\Rightarrow$ unbroken centre $\Rightarrow$ confining}
(Chatterjee 2021). The spectral gap is a \emph{sufficient-condition input} to the area law, not a
corollary of it. \LN
\end{keybox}
\src{\code{MU5-SIGMA-M.md} \S1.2 (quoting Chatterjee, \emph{CMP} 2021, \texttt{arXiv:2006.16229},
\texttt{doi:10.1007/s00220-021-04086-y}, verbatim: ``if \dots correlations decay exponentially \dots,
then center symmetry does not break'').}

\noindent
At the same time a widely repeated citation must be corrected: the abstract of Osterwalder--Seiler
1978 announces only a ``positive self-adjoint transfer matrix'', existence and analyticity of the
infinite-volume limit, and ``\textbf{we verify Wilson's confinement bound}'' --- \textbf{the abstract
contains no mass gap}. ``OS 1978 proved that the area law implies a gap'' has no checkable textual
basis. \LT
\src{\code{MU5-SIGMA-M.md} \S1.1 (checked verbatim against the INSPIRE/Elsevier abstract,
\texttt{doi:10.1016/0003-4916(78)90039-8}); the same file lists seven ``not found'' items in \S7 with
the query strings.}

\subsection{The area law carries its own gapless direction}

This is the hardest statement of the section. The area law's own correction term (effective string
theory, the Lüscher term)
\begin{equation}
  V(R)=\sigma R-\frac{\pi(D-2)}{12R}+O(R^{-3})
\end{equation}
shows that \emph{in the very sector that produces the area law the level spacing is
$O(1/R)\to0$} --- the string's oscillation modes. \textbf{The area law carries a gapless direction in
the $R\to\infty$ limit}, so $\sigma>0$ cannot by itself entail a bulk gap; the gap must come from
another sector (local gauge-invariant operators), and the area law is precisely silent about that
sector. \LT
\src{\code{MU5-SIGMA-M.md} \S2 and \S3.1; Lüscher--Symanzik--Weisz, \emph{NPB} \textbf{173} (1980)
365 (checked through secondary quotation; that file marks it \code{[B]}).}

\noindent
\textbf{The exact missing step (fixed in writing)}: from ``area law'' to ``exponential decay of
connected correlators of local gauge-invariant operators'' (item L4 of that file). It needs
interchange of limits, uniqueness of the vacuum and OS reconstruction, none dispensable; it does not
exist in the literature and cannot be assembled from L1--L3. \LM
\src{\code{MU5-SIGMA-M.md} \S2 table and the ``exact missing step'' paragraph.}

\noindent
The two sides are indeed mutually independent, and the other side is also proved: the
fundamental-representation Higgs region (Fradkin--Shenker type) has a mass gap while Wilson loops obey
a perimeter law (no area law). And $\sigma=m^{2}$ is false in every model --- for 4D
$\mathrm{SU}(3)$, $m^{2}/\sigma\approx15$; in 2D YM both scale like $g^{2}N$ and the ratio is a pure
number, \emph{not} an equality. \LT
\src{\code{MU5-SIGMA-M.md} \S3.2 and \S4.}

\subsection{Two things this section leaves open (stated honestly)}

\begin{itemize}
\item \textbf{The opposite claim $\sigma\propto\Lambda^{2}$.} \code{MU6-NUM.md} uses
      $\sqrt{\sigma}=484.96$ MeV and $\Lambda_{\mathrm{MS}}=263.04$ MeV to get
      $\sigma/\Lambda^{2}_{\mathrm{MS}}=3.399/2.967$, with $2g^{2}N$ run to $\Lambda_{\mathrm{MS}}$
      giving $9.363$ and ratios $0.567/0.363$, i.e.\ $O(1)$, and local exponents
      $n=2.63\,(\mathrm{SU}(3))/2.31\,(\mathrm{SU}(2))$, concluding $\sigma\propto\Lambda^{2}$. This
      contradicts $n=0$ from \code{mu-4}. \LM
\item \textbf{A normalization error found in the same round}: the anchor $a^{-1}(\beta=6)$ of $2.91$
      GeV is high by $1.30\times$; the correct value is $2.243$ GeV (AT2021 Tab.13). The direction of
      the conclusion is unchanged, but every number that used $2.91$ must be divided by $1.30$. \LE
\end{itemize}
\src{\code{LOOP-RUNBOOK.md} \S11.37 (lines 1804--1823); \code{MU6-NUM.md}. This document's handling:
both sides are displayed and not adjudicated; but \emph{the string-tension side as a route to $m>0$}
is closed by the four steps above, independently of this conflict --- the closure rests on
prescription dependence and on the direction of the proved arrow, neither of which depends on the
value of $n$.}

% =====================================================================
%  §6  Structural closure of the loop algebra
% =====================================================================
\section{Structural closure of the loop algebra: no sector is both faithful and non-degenerate}
\label{sec:loop}

The loop-equation route had two hopes: (a) use scale covariance to pin the string tension as
$\propto\Lambda^{2}$; (b) find on non-planar loops a genuine test of the area law. Both closed on the
same day, but in different ways: the first turned out to be a \emph{weak} result (it excludes
homogeneous monomials only), the second turned out to be \emph{impossible} (no sector is
simultaneously faithful and non-degenerate).

\subsection{$\kappa_W=0$: a true theorem, but it excludes only homogeneous monomials}

Apply the scaling $x\to sx$, $C_R\to C_{sR}$ with $\theta$ fixed to \eqref{eq:MM}: the left side
$\partial_\mu\delta/\delta\sigma_{\mu\nu}$ carries weight $s^{-3}$, and so does the right side
$\oint\dd y\,\delta^{4}$. Writing $W(sC)=s^{\kappa_W}W(C)$ and $G=W\cdot W$ (so $\kappa=2\kappa_W$),
matching gives
\begin{equation}
  \kappa_W-3=2\kappa_W-3 \;\Longrightarrow\; \kappa_W=0 ,
\end{equation}
and the solution is unique (it is pinned by the power difference between $O(W^{2})$ and $O(W)$, not
by the area law or by positivity). A second, independent derivation: a zero-length arc gives
$G(R,0)=W(R)$, hence $\kappa_G=\kappa_W$, while the product gives $\kappa_G=2\kappa_W$; the same
solution. \LT
\src{\code{RV-3-COVARIANCE.md} \S1--\S2; \code{LOOP-RUNBOOK.md} \S11.19, \S11.28.}

\noindent
\textbf{Its entire non-empty content is one statement: zero anomalous dimension} --- $W$ carries no
$R^{-p}$ multiplicative prefactor. It excludes every pure power decay $W\sim R^{-p}$ with $p\neq0$,
and it excludes schemes in which $\kappa_W$ is tuned as a free parameter. \LT
\src{\code{KW1-CONS.md} \S1--\S2, \S5.1.}

\begin{keybox}
\textbf{Stated plainly: this is a weak result.} For $F(x)=e^{-cx^{\alpha}}$ one has
$\kappa_W(s)=-\alpha c(sx)^{\alpha}$, which is \emph{not constant} (unless $\alpha=0$ or $c=0$), so
``$\kappa_W=0$'' \textbf{does not apply} to it --- not ``it is allowed'', but a \emph{type error}.
Reductio: reading $\kappa_W=0$ as a pointwise constraint on profiles would exclude
$e^{-\sigma R^{2}}$; but \eqref{eq:sigma-mu} of \S\ref{sec:tension} already established that the area
law \emph{is} a solution $\Rightarrow$ that reading is wrong $\Rightarrow$ \textbf{$\alpha$ is entirely
free}. \LT
\end{keybox}
\src{\code{KW1-CONS.md} \S1.2(ii), \S2.2; numerical check in \S1.3 ($s=0.5,1,2,8$ columns, with
$e^{-x},e^{-x^{2}},e^{-x^{1/2}},e^{-x^{3}}$ all non-$s$-constant and $x^{-2}$, constant both
$s$-constant).}

\noindent
\textbf{The real source of $\alpha=2$ is elsewhere}: in the additive / planar sector \eqref{eq:MM}
reduces to $G_R=-(4\pi g^{2}N/\varepsilon^{2})RG$, i.e.\ the source has $R$-power $q=1$; with
$\alpha=q+1$ this gives $\alpha=2$ and $\sigma=2g^{2}N/\varepsilon^{2}$, valid in that sector only and
independent of $\kappa_W$. \LT
\src{\code{KW1-CONS.md} \S4.2 (including the four-point check of $G_R/G$ at $R=0.5,1,2,4$);
\code{RV-1-EQUATION.md} \S8.3(i).}

\noindent
\textbf{An attribution error (erratum level)}: crediting $\alpha=2$ to $\kappa_W$ is wrong;
$\kappa_W=0$ has \textbf{zero content} for $\sigma$, for the spectrum and for $m_{0^{++}}$, hence it
\textbf{cannot be used to prove a gap}. \LN
\src{\code{KW1-CONS.md} \S3.4 attribution table (six independent rows, the first five with different
owners), \S5.3(4), \S5.4; \code{LOOP-RUNBOOK.md} \S11.36 (lines 1780--1802, ``I overrated
$\kappa_W=0$'').}

\subsection{The area law holds only in the degenerate sector}

On the planar circle family $A_1+A_2\equiv\pi R^{2}$, so $\theta$ always drops out (step 1 of
\S\ref{sec:tension}). The mode language makes this precise: in the basis
$\varphi_k=\cos(2\pi k\theta)$ the $k=0$ equation involves $G_0$ only (\emph{closure}), but closure is
\emph{not} degeneracy: if $S_k=0$ and there is no resonance then $G_k=0$ ($\theta$-independence), which
contradicts the physical $G=W_1W_2$ (generically non-multiplicative) $\Rightarrow$ the physical $k\ge1$
components must be carried by the contact source. The checkable criterion is
\begin{equation}
  A_1(R):=2\int_0^{1}G(R,\theta)\cos(2\pi\theta)\,\dd\theta ,
  \qquad A_1\neq0 \Rightarrow \text{the two variables cannot be dispensed with}.
\end{equation}
\LT
\src{\code{RV-3-COVARIANCE.md} \S3, \S5 (with the numerical check that the Fourier coefficients of
$\theta(1-\theta)$ reproduce $-1/(\pi^{2}k^{2})$ digit by digit); \code{RV-4-LINEAR.md} for the
correction that the second variable brings a $\theta$-direction second-order term plus a resonance
condition, not a nonlinearity.}

\noindent
\textbf{So the planar family cannot refute the area law}: it tests the coefficient, not the functional
form. Testing the functional form requires an \emph{non-additive} splitting, i.e.\ $A_1+A_2\neq A$ ---
a non-planar loop or a self-intersecting one. \LT
\src{\code{RV-2-AREALAW.md} \S6.}

\subsection{The non-degenerate sector: the principal symbol is never elliptic}

Take the two-variable reduction on a figure-eight. First fix the geometry: two coplanar circles with
\emph{exactly one} common point $\iff$ the radical axis is tangent to the circles $\iff$ the circles
are \textbf{tangent} (parallel tangents, oriented angle $\pi$). ``Coplanar $+$ one common point $+$
angle $\pi/2$'' does not exist geometrically (a transverse intersection always gives two points). \LT
\src{\code{FIG8.md} \S1.1 (the radical-axis argument).}

Tangency kills two things at once: \LT
\begin{itemize}
\item the right-hand $\dd y_\nu$ weight satisfies $t_1+t_2=0$ (because of $2\cos(\varphi/2)$ with
      $\varphi=\pi$) $\Rightarrow$ \textbf{RHS $\equiv0$};
\item $\rho_2=\partial x_*/\partial R_2=0$ $\Rightarrow$ the $\partial^{2}/\partial R_2^{2}$ slot is
      \textbf{permanently empty}.
\end{itemize}
Hence the tangent figure-eight is not a two-variable reduction but ``one variable $+$ one empty slot''
$+$ zero right-hand side: a first-order homogeneous constraint. \LN
\src{\code{FIG8.md} \S2.2--\S2.3, \S7 item 1.}

\noindent
A genuinely transverse figure-eight (non-circular, or two skew circles) gives nonzero $\rho_b$ and
$\varphi\neq\pi$. But its principal symbol matrix
\begin{equation}
  \begin{pmatrix} a_{11} & a_{12}/2\\ a_{12}/2 & a_{22}\end{pmatrix},
  \qquad
  \det \;=\; -\frac{(\rho_1R_1-\rho_2R_2)^{2}}{16\pi^{2}R_1^{2}R_2^{2}} \;\le\; 0 ,
\end{equation}
with equality only on one hypersurface. That is: \textbf{at every non-degenerate point the operator is
hyperbolic, never elliptic}. \LT
\src{\code{FIG8.md} \S0.5 item 3 (with the five-point numerical check of the $\det$ formula, maximum
deviation $3.5\times10^{-18}$).}

\begin{keybox}
\textbf{Immune-graph family 57 (the third closure of this document).} Maximum-principle arguments
(maximum principle, Harnack, comparison theorems, monotonicity) \emph{all require ellipticity}. A
hyperbolic operator has no maximum principle $\Rightarrow$ any proposal to ``get a lower bound on loop
space from monotonicity / a maximum principle'' is \textbf{structurally inexpressible} in loop-equation
coordinates --- it is not ``not yet done''. \LN
\end{keybox}
\src{\code{LOOP-RUNBOOK.md} \S11.23 (lines 1546--1551); the family-57 entry of \code{H1B-FIVE.md}
\S C2; \code{FIG8.md} \S0.5 item 3 for the four-step chain (sign of $f''$ $\Rightarrow$ monotonicity of
$f'$ $\Rightarrow$ comparison via $f^{2}$ homogeneity $\Rightarrow$ boundedness of $1/f$).}

\noindent
\textbf{The area law is not a solution in this sector, and the reason is dimensional.} Substituting
term by term gives
\begin{equation}
  \frac{\text{LHS}}{G}=4\pi\sigma^{2}(\rho_1R_1+\rho_2R_2),
  \qquad R\text{-power} = +1 ,
\end{equation}
so the residual is $O(\sigma^{2}R)\to\infty$ and there is no solution for any $\sigma\neq0$ (at
$\sigma=0$, $G\equiv1$ and only $b_1=0$ survives). Numerically ($\sigma=0.37$):
$4\pi\sigma^{2}R_1=1.7203361\,R_1$, five points agreeing to $10^{-5}$; sweeping $R_2$ from $0.7$ to
$10$ changes LHS/$G$ only through $R_1$ (the numerical trace of $a_{22}=0$); and
LHS/$G\div R_1$ at $R_1=1,2,4,8$ is constant to the $1.7206$ level $\Rightarrow$ the power is exactly
$+1$. \LT+\LE
\src{\code{FIG8.md} \S4.1--\S4.3; witness \code{w7/fig8_res.py}.}

\noindent
And this does not depend on any geometric factor: \textbf{any} second-order $L^{-3}$ reduction fails to
hold a Gaussian, because $[a_{ij}]=L^{-1}$ is fixed by $L^{-3}$ while
$\partial^{2}_RG/G=O(\sigma^{2}R^{2})$, so $a\cdot\partial^{2}G/G=O(\sigma^{2}R)$ (power $+1$), while
RHS/$G=O(R^{0})$. Adding variables, changing the intersection angle, or going non-planar only changes
the pure-number coefficients $\rho_b$, $2\cos(\varphi/2)$, never the $R$-power. \LT
\src{\code{FIG8.md} \S4.2 (``fix the slots first, fill in the geometry last''); this is precisely the
antidote to the earlier $c_3$ mishap.}

\subsection{One sector, two requirements, not jointly satisfiable}

\begin{table}[h]
\centering\footnotesize
\setlength{\tabcolsep}{6pt}
\renewcommand{\arraystretch}{1.12}
\begin{tabularx}{\textwidth}{@{}l X X l@{}}
\toprule
\textbf{Sector} & \textbf{Faithful (area law is a solution)} & \textbf{Non-degenerate} & \textbf{Conclusion}\\
\midrule
circle (\code{rv-1}/\code{rv-2}) & \textbf{yes} (residual identically $0$, $1.62\times10^{14}\to0$) & \textbf{no}
($A_1+A_2\equiv\pi R^{2}$ $\Rightarrow$ pointwise algebraic; $G_\theta\equiv0$ forced) & tests the coefficient only\\
figure-eight (\code{fig8}) & \textbf{no} (residual $4\pi\sigma^{2}(\rho_1R_1+\rho_2R_2)$, power $+1$) & \textbf{yes}
($\det\le0$, non-congruent, $H2$ fails) & structure without faithfulness\\
\bottomrule
\end{tabularx}
\caption{The final death cause of the loop-algebra line: no sector satisfies both requirements.}
\label{tab:sectors}
\end{table}

\src{\code{LOOP-RUNBOOK.md} \S11.22 (lines 1536--1544); \code{FIG8.md} \S0.5 items 1--3, \S4, \S7.}

\begin{keybox}
\textbf{Structural conclusion (the main result of this section).} The loop algebra has \textbf{no sector
that is simultaneously faithful and non-degenerate}: the circle sector yields
$\sigma=2g^{2}N\mu^{2}$ but degenerates into an algebraic relation (it tests the coefficient only);
the figure-eight has genuine structure (non-multiplicative, $k\ge1$ nonzero, $\det\le0$) but the area
law is not its solution (residual power $+1$). \LN
\end{keybox}
\src{\code{LOOP-RUNBOOK.md} \S11.22, last paragraph (verbatim: ``the final death cause of the string
tension line (SP-F)'').}

\subsection{Incidentally: the linear layer does not force a gap either}

The linearized layer forces no decay at all: the surviving solutions are a power law
$R^{-0.679}$ and $G\to c\neq0$. But $p=-0.679\neq0$ is \emph{forbidden} by $\kappa_W=0$, so that
power-law solution is an artefact of the bookkeeping, not a solution; what remains, $G\to c\neq0$, is
not decay. Hence the linear layer has zero content. \LN
\src{\code{RV-4-LINEAR.md}, as reported in \code{LOOP-RUNBOOK.md} \S11.16; the sentence ``$p=-0.679$ is
forbidden'' is at \code{LOOP-RUNBOOK.md} line 1647 and \S11.28.}

% =====================================================================
%  §7  Literature verdict
% =====================================================================
\section{Literature verdict: the shape is known, the theorems are zero}
\label{sec:lit}

This section separates two things on the literature side: \textbf{the shape} ($m=C\Lambda$) has
numerical support and two-dimensional precedent; \textbf{$a$-uniform theorems} are zero. Both hold at
once, and that is precisely the external version of the verdict of \S\ref{sec:position}.

\subsection{The shape: $m/\Lm=6.3(3)$ at the continuum limit, three independent routes agreeing}

\begin{table}[h]
\centering\footnotesize
\setlength{\tabcolsep}{5pt}
\renewcommand{\arraystretch}{1.12}
\begin{tabularx}{\textwidth}{@{}l X l@{}}
\toprule
\textbf{Route} & \textbf{Data} & \textbf{Result}\\
\midrule
$m/\sqrt{\sigma}\div\Lm/\sqrt{\sigma}$ & AT2021 (\texttt{arXiv:2106.00364}, JHEP12(2021)082) Tab.34 $\div$ Tab.15:
$\mathrm{SU}(3)$ $3.405(21)/0.5424(13)$; $\mathrm{SU}(2)$ $3.781(23)/0.5806(21)$ & $6.28(2)(21)$ / $6.51(2)(64)$\\
$r_0$ route ($r_0\Lm=0.629(4)$) & MP1999 \texttt{hep-lat/9901004} \S V: $r_0m=4.21(12)$; Chen2006 \texttt{hep-lat/0510074}
Tab.XXI: $4.16(11)(4)$ & $6.69(20)(23)$ / $6.61(18)(24)$\\
GeV division & AT2021 Tab.37: $1.653(26)$ GeV $\div$ eq.22: $0.263(4)(9)$ GeV & $6.28(14)(22)$\\
\bottomrule
\end{tabularx}
\caption{$m_{0^{++}}/\Lm^{(N_f=0)}=6.3(3)$: three mutually independent routes, with $\mathrm{SU}(2)$ and
$\mathrm{SU}(3)$ agreeing.}
\label{tab:c}
\end{table}

\noindent
Together with the uniqueness \eqref{eq:unique} of \S\ref{sec:position} the target is now: \textbf{prove
$m\ge c\Lambda$ with $c>0$ a pure number}; the value is known, $C=6.3(3)$ (about $5\%$ precision). What
is missing is the proof, not the number. \LE
\src{\code{L2-GLUEBALL-LAMBDA.md}, as reported in \code{LOOP-RUNBOOK.md} \S10.16--\S10.17 and
\code{LIT-VERDICT.md} \S6; the scheme ratio $\Lm/\Lambda_I=2.633$ is from \code{L2-GLUEBALL-LAMBDA.md}.}

\subsection{$a$-uniform theorems: zero (and not ``not found'').}

\begin{itemize}
\item \textbf{The only rigorous mass gap on the lattice}: Kupiainen, \emph{CMP} \textbf{73} (1980) 273
      (\texttt{doi:10.1007/BF01197703}) proves $m>0$ and a $1/n$ expansion, \textbf{with no shape}. \LT
\item \textbf{Constructive 4D gauge theory}: Magnen--Rivasseau--Sénéor, \emph{CMP} \textbf{155} (1993),
      title verbatim ``\emph{with an infrared cutoff}'' --- the infrared is cut off, no $m$; the Balaban
      series (\emph{CMP} 102/109/116) proves UV stability and a small-field effective action, no $m$. \LT
\item \textbf{The only $m$ bound is outside that line}, Osterwalder--Seiler 1978: $ma\ge c(\beta)$, so
      $a\to0$ gives $m\ge c(\beta)/a\to\infty$ --- \textbf{the direction is reversed}. \LT
\item \textbf{Precedents for the shape exist only in 2D and in SUSY}: seven of them (O(3)/O(4),
      O($N$) $\sigma$ models, SUSY $\mathbb{CP}^{N-1}$, SUSY O($N$), etc.), all relying on
      integrability or SUSY; 4D has neither. \LT
\end{itemize}
\src{\code{LIT-VERDICT.md} \S1--\S4 (every DOI checked); \code{L1-UNIFORM-A.md},
\code{L3-TRANSMUTATION.md}, \code{L5-CLASSICAL.md}.}

\begin{table}[h]
\centering\footnotesize
\setlength{\tabcolsep}{5pt}
\renewcommand{\arraystretch}{1.12}
\begin{tabularx}{\textwidth}{@{}>{\raggedright\arraybackslash}p{3.5cm} >{\raggedright\arraybackslash}p{3.6cm} X@{}}
\toprule
\textbf{Source} & \textbf{Verbatim} & \textbf{Verdict}\\
\midrule
\texttt{arXiv:2606.19362}\newline(Faizal--Shabir, \emph{Fortschr.\ Phys.} \textbf{74} (2026) e70097)
 & ``transfer operator with a uniform gap \dots\ endure in the continuum''
 & \textbf{downgraded to a claim}: no constant, no display number; the core is \code{Hypothesis 1};
 uniform-in-$a$ comes from a \emph{postulated admissible class}; contraction to the Gaussian fixed
 point; no numerical bound on $\Delta_{\mathrm{phys}}$\\
\texttt{arXiv:2506.00284}\newline(Jacobsen) & ``no remaining gap'' & \textbf{excluded}: arXiv marks it
 withdrawn (``does not meet arXiv's research content quality standards'')\\
\texttt{arXiv:2307.00788}\newline(Lim) & gap $=$ Casimir label & dimensionless, no scaling, not a lattice
 $a\to0$ limit\\
\bottomrule
\end{tabularx}
\caption{The three preprints claiming to settle 4D: all die at the same step.}
\label{tab:preprints}
\end{table}

\noindent
The furthest point reached by the other nine is ``fixed-lattice-spacing strong coupling / 't Hooft
region'', and the field's own verdict is already in print: \texttt{arXiv:2611.11534} (2026-11-08)
verbatim ``\emph{A nontrivial four-dimensional spatial-cutoff construction and a positive spectral
bound along it are not established.}''; Zenodo 23783328 verbatim ``\emph{the composition across
scales, the large fields and the crossover remain open. The cube estimate and the existence of the
continuum limit are not proved, so the continuum Yang--Mills existence and mass-gap problem remains
open.}'' \LT
\src{\code{LIT-VERDICT.md} \S4--\S5 and \S7; \code{L4-PREPRINTS.md} (twelve IDs, each opened).}

\begin{keybox}
\textbf{The common gap of the three claims is the object of \S\ref{sec:five}.} Each of them, in
transporting ``strong-coupling gap $\to$ weak-coupling continuum window'', turns a hypothesis into a
theorem; and that step is exactly the rate uniformisation of $C5$, i.e.\ $d5/(CP)=p=1$. This is the
strongest external confirmation in this document: \textbf{every attempt fails at the same step, and we
know the step's name.} \LT
\end{keybox}
\src{\code{LOOP-RUNBOOK.md} \S10.18 (lines 1006--1019); \code{LIT-VERDICT.md} \S7, conclusion 4.}

\subsection{A premise that must be attached to $6.3(3)$: it is conditional on the scaling form}

The value $6.3(3)$ is obtained by extrapolating $m/\sqrt{\sigma}$ and $r_0m$ measured at
\emph{accessible} $\beta$ to the continuum limit, so it depends on the scaling form $m=C\Lambda$ itself.
At accessible $\beta$ the measured logarithmic slope is only $0.53\,(\mathrm{SU}(3))$ and
$0.21\,(\mathrm{SU}(2))$ of the asymptotic value, and for $\mathrm{SU}(2)$ the $N_\tau=4$ data even
have the opposite sign; reaching $90\%$ of the asymptotic slope needs $\beta_{90}\approx27$ (where
$\xi/a\sim10^{11}$), which is inaccessible. \LE
\src{\code{C2-MEASURED-SLOPE.md}, as reported in \code{LOOP-RUNBOOK.md} \S10.20 (including the
extrapolation formula for $\beta_{90}$ and the $1\sigma$ interval $[6.3,58]$).}

\begin{keybox}
\textbf{Hence $6.3(3)$ is ``compatible with the scaling form'', not ``established by it''.} The
numerical route \textbf{cannot independently confirm} the shape --- it takes the very thing to be proved
as the hypothesis of its extrapolation. \LN
\end{keybox}
\src{\code{LOOP-RUNBOOK.md} \S10.20, the honest-premise paragraph (lines 1064--1067).}

\noindent
This is the same statement as \S\ref{sec:position}, seen from the literature side: the shape is
numerically confirmed to $5\%$ precision, while the route to it presupposes $m>0$. That is what the
title means by ``a consequence, not a route''.


% =====================================================================
%  §8  Errata
% =====================================================================
\section{The day's errata: 26 normalization, coordinate, convention and enclosure errors}
\label{sec:errata}

The campaign caught 26 instances of ``a quantity placed where it does not belong'' in a single day.
They are not a collection of typos but twenty-six instances of one family: \textbf{putting a quantity
where it does not belong} (the lesson of \code{LOOP-RUNBOOK.md} \S10.17, quoted verbatim). Several of
them decided whether a verdict could be pronounced at all, so each is written as
``wrong $\to$ right $\to$ source'', with a marker for the ones that changed a verdict.

\vspace{2pt}
{\footnotesize
\setlength{\tabcolsep}{3pt}
\renewcommand{\arraystretch}{1.05}
\begin{longtable}{@{}>{\raggedleft\arraybackslash}p{0.45cm} >{\raggedright\arraybackslash}p{1.95cm} >{\raggedright\arraybackslash}p{4.15cm} >{\raggedright\arraybackslash}p{4.9cm} >{\raggedright\arraybackslash}p{3.65cm}@{}}
\toprule
\textbf{\#} & \textbf{Kind} & \textbf{Wrong} & \textbf{Right} & \textbf{Source}\\
\midrule
\endfirsthead
\toprule
\textbf{\#} & \textbf{Kind} & \textbf{Wrong} & \textbf{Right} & \textbf{Source}\\
\midrule
\endhead
\bottomrule
\endfoot
1 & normalization & the number $c=12$--$18$ for $m/\Lambda$ & $c=6.3(3)$; main cause: the corpus value
$ma=1.073(88)$ was taken as the physical $0^{++}$ (too large by $1.89\times$); secondary cause:
$\Lambda$ taken as $200$--$240$ MeV instead of $263(4)(9)$ & \code{L2-GLUEBALL-LAMBDA.md};
\code{LOOP-RUNBOOK.md} \S10.16 \textbf{(verdict changed)}\\
2 & convention & $m/\Lambda=21$--$25$ (\code{a10}) & withdrawn: two independent errors $\times1.889$ and
$\times0.558$, net $3.38$ (matching the measured ratio) & \code{LOOP-RUNBOOK.md} \S10.17
\textbf{(verdict changed)}\\
3 & domain & using two-loop running to invert $\Lambda_{\mathrm{MS}}$ at $g^{2}=1$
($\beta=6$, $\mathrm{SU}(3)$) & invalid: gives $147$ MeV against the physical $263$ MeV (off by
$1.79\times$); the two-loop formula is right, the domain is wrong & \code{LOOP-RUNBOOK.md} \S10.17\\
4 & coordinate & $\lambda\sim1$ (Balaban's stopping point, reading $g^{2}$) and $\tH=1$ (lower edge of the
crossover, reading $\beta N$) treated as the same $1$ & same name, different coordinates, differing by
exactly $2N^{2}$ ($\mathrm{SU}(3)$: $18$; $\mathrm{SU}(2)$: $8$); one must cross $2.47\times$
(physical point) / $18\times$ (lower edge), not $7.3\times$ & \code{BC-CHAIN.md}:90;
\code{LOOP-RUNBOOK.md} \S9.8\\
5 & numerical & $\Gpp(\beta=2.5)=+0.5427$ & $+0.135995$: the second-order term of $\sin(\rho/2)u$, namely
$s\cos(r/2)/(2r)$, was double counted; $u''=0$ and $\rho'=0$ make it vanish. The test is the archived
$E[H]\equiv0$ & \code{H1A-EXACT.md} \S2.4\\
6 & normalization & $Q(\infty)=0.1099$ & $0.2667$ (since $0.1099=0.32957/\lambda_1$, a double error from
dividing the 13-point trapezoid by $\lambda_1$); plus $Q\propto\beta^{1.14}$, $p=0$, and the threshold
$e^{-Q}\approx0.77$ instead of $0.90$ & \code{A3-Q-SCALING.md}; \code{LOOP-RUNBOOK.md} \S10.8
\textbf{(verdict changed)}\\
7 & norm & reading $L^{4}$ as ``curvature blow-up'' & $L^{4}$ is the matrix dimension: bounding the
operator norm by the trace loses exactly $\dim=3E=12L^{4}$; measured $\mathrm{tr}/$row-sum
$=126/1642/10364$ with row sums $=3.51/4.31/3.40=O(1)$ & \code{DENSITY-D.md}; \code{LOOP-RUNBOOK.md} \S9.4\\
8 & sup/inf & using $d=9.375$ as a $\sup$ & it is a \emph{lower} bound: the correct interval is
$d\in[3.62,6]\beta^{2}$ ($\beta=2.5$: per plaquette $[22.64,37.5]$); exponents and verdicts unchanged,
the constant was underestimated by $\ge2.2\times$ & \code{R3-L4.md}; \code{LOOP-RUNBOOK.md} \S9.9\\
9 & normalization & $\operatorname{Hess}\succeq-(4.8/N)\beta$; later $-(2.9/N)\beta$ & the first is the
graph-cardinality bound without symmetrization ($\sup\lambda_{\max}=0.752$--$0.763$); the second
symmetrizes ($\tfrac12[\lambda,\lambda]$) but still sums per plaquette (missing cross-link blocks, a
factor $8/3$) $\Rightarrow$ the truth is $-8\beta/N$; consequently $\dot\lambda_0$ changes from $-3.51$
to $-6.29$, threshold $\beta<N^{2}/64$ & \code{AT-STAR.md} E3; \code{MC-METRIC.md} E1;
\code{LOOP-RUNBOOK.md} \S9.7\\
10 & coordinate & $120\times$ and $13.3\times$ (distance to the SZZ radius) & $\beta_W=N^{2}\beta_A$;
correct values $27.6\times(\mathrm{SU}(2))$ / $30.4\times(\mathrm{SU}(3))$, i.e.\ about $30\times$.
Consequently $1/32\to8.9\times$ should read $19$--$20\times$, and $\tH<N/48$, $\lambda\approx7.3$ should
read $\lambda>96$, $\lambda\approx3.2$ & \code{BETA-CONVENTION.md}; \code{LOOP-RUNBOOK.md} \S9.8
\textbf{(verdict changed)}\\
11 & normalization & using $u_{P3}=c_1/c_0$ as the per-tile weight & $u_{\mathrm{tile}}=I_2/I_1=u_{P3}/2$;
hence $m_1a(2.5)=+2.715438$, not $-0.057151$ & \code{R4-PERCENT.md} E7; \code{LOOP-RUNBOOK.md} \S9.3\\
12 & convention & $\beta_{\mathrm{tube}}=2.446918$ and ``off by $2.12\%$'' & withdrawn: a threefold
misalignment ($u_{P3}$ normalization $\times$ $|X|=1$ $\times$ threshold $1\neq a$); the direction is not
``$2\%$ off'' but \textbf{$111\%$ over}; correct $\beta_c^{KP}\le1.182227$ (per site), optimal $a=1/6$.
Also the tabulated $\varphi(\text{cube})$ $2.03\times10^{-3}$ should be $1.916\times10^{-3}$ &
\code{R4-PERCENT.md}; \code{LOOP-RUNBOOK.md} \S9.3; \code{AUDIT-R5B.md} \textbf{(verdict changed)}\\
13 & combinatorics & plaquette adjacency number $\Delta=12$ (a 3D value) & $\Delta=20$ ($4(2d-3)$,
verified by brute force for $d=2,3,4,5\to4,12,20,28$); the $\kappa$ upper bound moves from $48$ to $76$,
so $\kappa\le4.3307$ and the interval is $[1.4012,4.3307]$; also $\inf_A a_A^{1/A}$ was measured to be
increasing, so it should read $\lim$ & \code{Q5-KAPPA.md}; \code{LOOP-RUNBOOK.md} \S11.40
\textbf{(verdict changed)}\\
14 & dimension & $\mu^{2}=\delta^{3}_{\perp}(0)$ ($L^{-3}$) read directly as $\mu^{2}$ ($L^{-2}$) &
$\mu^{2}=\varepsilon\,\delta^{3}_{\perp}(0)=1/\varepsilon^{2}$; exactly one power of $a$ is missing (the
curve delta normalization $\delta_C=a^{-1}$), and it is not a convention that may be dropped &
\code{MU-3-RENORM.md} \S1; \code{MU-4-LATTICE.md} \S1.3, \S4\\
15 & sign & \code{(8.4)} giving $G_R=4\pi g^{2}NR\,G/\varepsilon^{2}$ coexisting with
$\sigma=+2g^{2}N/\varepsilon^{2}$ from \code{(8.5)} & they differ by a sign: either \code{(8.4)} needs a
minus on the right, or $(\star)$ does; the sign of $(\star)$ is not fixed & \code{MU-3-RENORM.md} \S1.2\\
16 & dimension & the one-dimensional reduction passing $f''$ off as $\partial_\mu\delta/\delta\sigma$
($L^{-2}$) & the operator is $L^{-3}$ as a whole; the power-law exponent moves from $p=-1$ to $p=0$
(which voids the ``rigid exact solution'' of $c_3$) & \code{SP-A-SPDE.md} \S0; \code{RV-3-COVARIANCE.md}
\S2; \code{LOOP-RUNBOOK.md} \S11.3\\
17 & convention & $0.2688$ taken as $\Gpp$ & it is $1/\operatorname{Var}(a)$ ($\operatorname{Var}=3.719824$);
$\Gpp$ at $\beta=0$ is $1/(2\,\mathrm{Cin}\,2\pi)=0.20511530$. The multiplicative reading
$D/S:=1-\Gpp\operatorname{Var}(a)$ gives $0.2370$ at $\beta=0$ while the true defect ratio is $0$ &
\code{DS-RATIO.md} (line 209)\\
18 & dimension & $(\star_1')$ right side $c_1f''+c_3f'/R=c_2g^{2}Nf^{2}/R^{3}$ & it should be
$c_2g^{2}Nf^{2}/R^{2}$ (one $R$ too many); this error alone voids the ``rigid exact solution'', whose
amplitude is an artefact of the wrong scaling & \code{C3-COEFF.md} \S1\\
19 & geometry & taking the intersection angle of the figure-eight to be $\pi/2$ & coplanarity plus a
single common point forces tangency, with oriented angle $\pi$; $b_0=2\cos(\varphi/2)=0$ $\Rightarrow$
RHS $\equiv0$, and $\rho_2=0$ (the $\partial^{2}/\partial R_2^{2}$ slot is permanently empty) &
\code{FIG8.md} \S1.1, \S1.4, \S2.3 \textbf{(verdict changed)}\\
20 & normalization & \code{na2}'s chain averages ``decreasing in $L$'' & $X_{na2}=X_{\text{true}}/L^{2}$
(too small by $L^{2}$); after rescaling $|X|\propto L$ and $X/L\to-0.010$ ($\beta=0.5$); the conclusion
that all $15/15$ are negative is unchanged & \code{CT-CHAIN.md} \S2.1; \code{NA2-NONDEGENERATE.md}\\
21 & citation & ``Osterwalder--Seiler 1978 proved area law $\Rightarrow$ gap'' & the abstract proves only
Wilson's confinement bound and a positive self-adjoint transfer matrix, and \textbf{contains no mass gap};
the proved implication runs $m>0\Rightarrow$ confinement (Chatterjee 2021) & \code{MU5-SIGMA-M.md}
\S1.1--\S1.2 \textbf{(verdict changed)}\\
22 & normalization & $\tH<1/24$ (as the SZZ-provable region) & no such display exists in the source; the
correct reading is $\tH<N/48$ (i.e.\ $\lambda=g^{2}N>96$), and the old $\lambda\approx7.3$ is wrong &
\code{LIT-CHECK.md} Q3$'$; \code{AUDIT-R1.md} \S A2; \code{LEDGER.md} \#7\\
23 & bookkeeping & $\sigma a^{2}=\ln(1/u)$; $\sigma>\kappa/a^{2}$ & the $-\kappa$ is missing: correctly
$c(\beta)=-\log u=\kappa+\sigma a^{2}/\pi$; the version with $\pi$ is wrong & \code{Q4-LOGCONVEX.md};
\code{LOOP-RUNBOOK.md} \S11.38(4)\\
24 & normalization & $a^{-1}(\beta=6)=2.91$ GeV (used as the anchor of \code{mu-3}/\code{mu-4}) & $2.243$
GeV (AT2021 Tab.13); high by $1.30\times$; the direction of the conclusion is unchanged but every number
that used $2.91$ must be divided by $1.30$ & \code{MU6-NUM.md}; \code{LOOP-RUNBOOK.md} \S11.37\\
25 & enclosure & $\mathrm{gap}(0)=0.729924127576750\ldots$ with an interval of width $1.4\times10^{-10}$,
quoted as a rigorous enclosure & $0.729924128620548449$ with width $\pm10^{-55}$: the four molecular
integrals and their Cauchy tails are valid enclosures, but the quotients $E[g]$, $E[\Phi'']$ were printed
with collapsed endpoints and the tails never propagated through the division, so the interval did
\emph{not} contain the truth (it lay $9.05\times10^{-10}$ outside, $6.5$ times its own width); the same
defect exceeded the $\theta=\pi/2$ width by $50\times$. Found by an independently written interval class,
and flagged independently from a second pair of coordinates & \code{R10-AUDITCERT.md} \S0, \S3;
\code{R06-GAPMONO.md} \S4 \textbf{(enclosure corrected)}\\
26 & bookkeeping & $L=2$ described as a $1\times1\times1\times1$ torus with $16$ links and $2$
plaquettes & $4L^{4}=64$ links, $6L^{4}=96$ plaquettes, $16$ sites; the $16$-link object is the blocked
lattice $L_{c}=1$ obtained from $L=2$ by a factor-$2$ blocking. The dimension count $192-48-1=143$
follows the corrected bookkeeping & \code{R04-L2CERT.md} \S1.0, with \code{O24-L2-PUSH.md}:49 as the
independent confirmation; the text corrected is \code{Q11-CERT.md} \S6\\
\end{longtable}
}

\subsection*{How to use this table (three rules)}

\begin{enumerate}[label=(\roman*)]
\item The seven rows marked ``(verdict changed)'' (1, 2, 6, 10, 12, 19, 21) are the ones that really
      affect conclusions; row 25 changes no sign but withdraws an enclosure, which for a certificate is
      the same kind of damage. The rest affect only numbers or wording. Every row was checked line by
      line, and none of them changes the sign verdicts of \S\ref{sec:tension} or \S\ref{sec:loop}. \LT
\item \textbf{Rows 4 and 10 are two attacks of the same disease}: there are two $\beta$ coordinates
      ($\beta_W=N^{2}\beta_A$). Wherever a bare $1/48$, $1/24$ or $1/32$ appears, first ask ``which
      $\beta$ is this?''. \LT
\item \textbf{Rows 14, 16 and 18 are another form of the same disease}: one power of $L$ is dropped or
      added. Wherever $\delta$, $\delta_C$, $\mu^{2}$ or $\partial_\mu\delta/\delta\sigma$ appears,
      write the slots down before computing. \LT
\end{enumerate}

\subsection*{One conclusion-level correction, made the same day}

The table above collects normalization, coordinate and convention errors only. One \emph{conclusion}
was also withdrawn on the same day and is recorded here rather than in the table: the inference
``$\beta_c(\mathrm{SU}(3))>\beta^{*}$ implies H1 fails before the $\mathrm{SU}(3)$ scaling window opens,
hence H1 is the wrong target for the Clay object'' assumed $\beta^{*}$ to be $N$-independent. The $N=3$
measurement of \eqref{eq:n3} shows $V_0(3)<0$, so $\beta^{*}(3)$ does not exist and H1 holds for every
$\beta>0$ at $N=3$ --- in particular on the whole $\mathrm{SU}(3)$ scaling region. The structural law of
\eqref{eq:n32} puts the zero at $\beta^{*}(3)\approx19.6$, i.e.\ $\beta^{*}/\beta_c=1.87\to3.45$; the
monotonicity of the measured $\delta$ only pushes it further out. \LE
\src{\code{LOOP-RUNBOOK.md} \S11.39 (lines 1879--1885) for the withdrawn conditional; \eqref{eq:n3} and
\S\ref{sec:failure} for its falsified antecedent. The measurement is a single-plaquette/patch Monte
Carlo, not the full theory.}

\vspace{3pt}
\noindent
\textbf{One lesson, quoted:} all twenty-six are the same error, ``putting a quantity where it does not
belong''. The operational form used from now on is: \emph{state the slot, then fill in the numbers}. \LT
\src{The lesson is verbatim from \code{LOOP-RUNBOOK.md} \S10.17; the ledger is \code{LEDGER.md} (50 main
rows, 10 contradictions, 11 quantities of undetermined dimension). This table collects only rows whose
``wrong $\to$ right'' can be checked character by character; the contradictions and the
undetermined-dimension items of \code{LEDGER.md} are not counted among the 26. Row 25 is the only row of
the 26 that concerns an \emph{enclosure} rather than a number, and it is recorded here rather than in the
certificate's own section because the failure mode --- a bound quoted wider than it was --- is the same
failure mode as the other twenty-five.}


% =====================================================================
%  Appendix : source index
% =====================================================================
\section*{Appendix: source index}
\addcontentsline{toc}{section}{Appendix: source index}
\label{app:sources}

Every assertion in this document lands in one of the files below. ``Checked verbatim'' means the
abstract or a sentence of the text was read in full; ``measured'' means the number was produced by a
script and archived. \textbf{Unless marked \code{docs/}, all file names below live in
\code{research/v3/species/}.}

\vspace{3pt}
{\footnotesize
\setlength{\tabcolsep}{4pt}
\renewcommand{\arraystretch}{1.05}
\begin{longtable}{@{}>{\raggedright\arraybackslash}p{4.6cm} >{\raggedright\arraybackslash}p{10.5cm}@{}}
\toprule
\textbf{File} & \textbf{What this document uses from it}\\
\midrule
\endfirsthead
\toprule
\textbf{File} & \textbf{What this document uses from it}\\
\midrule
\endhead
\bottomrule
\endfoot
\code{H1A-EXACT.md} & all of \S\ref{sec:scalar} and \S\ref{sec:failure}: the identity, the four models,
$\delta(\beta)$, $\beta^{*}$, the external validation, the self-correction\\
\code{BETASTAR-N.md} & $W_0\propto N$, $V_0\propto1/N^{2}$, $\beta^{*}(N)\propto N^{3/2}$, the
$\beta^{*}/\beta_c$ table ($4.30661/2.2986=1.87$, $19.6/5.6924=3.45$), $\Phi''(0)=N/12$, the $a\beta$
term (\S\ref{sec:failure})\\
\code{H1B-FIVE.md} & the five-condition verdicts, $C5\Rightarrow C2$, the accounting error, the single
target (\S\ref{sec:five})\\
\code{CONTINUUM-PLAN.md} & the slope identity and the original objective of the campaign (and its
refutation) (\S\ref{sec:position})\\
\code{LIT-VERDICT.md} & the literature verdict, including $6.3(3)$, ``$a$-uniform: zero'', the 2026
preprints (\S\ref{sec:lit})\\
\code{MU-1-PRESCR.md} & four prescriptions, four values of $\mu^{2}$ (prescription dependence)\\
\code{MU-3-RENORM.md} & $\mu^{2}$ missing one power of $L$; the $\dd\ln(\sigma a^{2})/\dd\ln a$
criterion; $\sigma=2g^{2}Nm^{2}$ as bookkeeping\\
\code{MU-4-LATTICE.md} & the lattice power table, $\sigma\propto1/a^{2}$, $R\propto L^{0.700}$,
$\sigma/\Lambda^{2}=351/629/982$\\
\code{MU5-SIGMA-M.md} & $\sigma>0\not\Rightarrow m>0$, Chatterjee 2021, the OS 1978 abstract, the
Lüscher term, item L4\\
\code{MU6-NUM.md} & the opposite claim $\sigma\propto\Lambda^{2}$ and the $1.30\times$ anchor error in
$a^{-1}$\\
\code{RV-1-EQUATION.md} & the conventions of $(\star)$, $\sigma=2g^{2}N/\varepsilon^{2}$, $q=1$ (planar
sector only)\\
\code{RV-2-AREALAW.md} & the area law as a solution, $\rho\equiv0.5$, $A_1+A_2\equiv\pi R^{2}$, the
voided $c_3$ verdict\\
\code{RV-3-COVARIANCE.md} & the weight $s^{-3}$, uniqueness of $\kappa_W=0$, $k=0$ closure, the
$A_1(R)$ criterion\\
\code{RV-4-LINEAR.md} & the linear layer forces no gap; $R^{-0.679}$ is forbidden\\
\code{KW1-CONS.md} & $\kappa_W=0$ excludes monomials only, $\alpha=q+1$, the attribution table, the
weakness of the result\\
\code{FIG8.md} & the tangency argument, $b_0=0$, $\rho_2=0$, $\det\le0$, residual power $+1$, the
dimension theorem\\
\code{LOOP-RUNBOOK.md} & immune families 46--57, all verdicts of \S9--\S11, $Q(\infty)$,
$8.9\times\to19$--$20\times$, $6.3(3)$\\
\code{LEDGER.md} & the ledger of load-bearing numbers (50 main rows, 10 contradictions, 11
undetermined-dimension items)\\
\code{BETA-CONVENTION.md} & $\beta_A$ versus $\beta_W=N^{2}\beta_A$, $30\times$, $19$--$20\times$\\
\code{CT-CHAIN.md}, \code{NA2-NONDEGENERATE.md} & the $L^{2}$ normalization erratum of \code{na2}\\
\code{Q4-LOGCONVEX.md}, \code{Q5-KAPPA.md} & the quantitative no-go
($c(2.5)=0.679<1.40$), $\Delta=20$, the interval for $\kappa$\\
\code{A3-Q-SCALING.md} & $Q(\infty)=0.2667$, $Q\propto\beta^{1.14}$\\
\code{C1-SLOPE-EQUIV.md} & the three death causes of the slope identity (tautology / false both ways /
presupposition)\\
\code{C2-MEASURED-SLOPE.md} & the measured slopes $0.53/0.21$, $\beta_{90}\approx27$\\
\code{C3-UNIQUENESS.md} & uniqueness of $m=C(N,\text{scheme})\Lambda$, $\theta=0$ by definition\\
\code{L2-GLUEBALL-LAMBDA.md} & $m/\Lm=6.3(3)$ along three routes, $\Lm/\Lambda_I=2.633$\\
\code{L4-PREPRINTS.md} & the sweep of twelve preprints and the common gap of the three claims\\
\code{AT-STAR.md}, \code{MC-METRIC.md}, \code{DENSITY-D.md}, \code{R3-L4.md} & errata rows 7--9
($L^{4}$, $d=9.375$, $-(4.8/N)\beta$)\\
\code{R4-PERCENT.md} & $u_{\mathrm{tile}}$, $m_1a(2.5)$, $\beta_c^{KP}\le1.182227$, $111\%$\\
\code{AUDIT-R1.md}, \code{AUDIT-R5B.md}, \code{LIT-CHECK.md} & $\tH<N/48$, the tabulated
$\varphi(\text{cube})$, the nonexistent $1/24$\\
\code{DS-RATIO.md} & $D/S=\beta^{2}/\beta^{*2}$, the convention behind $0.2688$\\
\code{SP-A-SPDE.md}, \code{C3-COEFF.md}, \code{SP-F-LOOPALG.md} & the $L^{-3}$ slot erratum, $f^{2}/R^{2}$,
the one-dimensional failure\\
\code{docs/paper-v2/} & sources of definitions and display numbers: \code{ch03-existence.tex}
(dimensions), \code{ch05-theorems.tex} (\code{eq:split}, \code{eq:wgeo}), \code{ch08-immune.tex} ($(M)$),
\code{ch13-machine.tex} (five conditions, ceilings, \code{untouched}), \code{ch15-uvface.tex} (the
$\beta=0$ anchor, the periodic kernel), \code{ch19-h1-status.tex} (the certificate ledger: the three
certified values, the eight lemmas with their scopes, the two window grades, the lattice ledger and the
twelve roads), \code{app.tex} (\code{app:sharp}, \code{app:erratum-radius})\\
\midrule
\multicolumn{2}{@{}l}{\emph{Added by the certificate wave (all read in full unless marked)}}\\
\midrule
\code{Q11-CERT.md}, \code{R03-CERT-THETA.md}, \code{R10-AUDITCERT.md}, \code{R06-GAPMONO.md} & the
strict certificate of \S\ref{sec:scalar-cert}: the three values and their enclosures, the $800$-bit
pipeline with exact rationals, the two bugs found by the $\theta=\pi$ regression, the independent audit
that corrected the $\theta=0$ enclosure, the exponential-family rewriting, the five-term derivative of the
gap, and the $501$-point monotonicity with its two complement routes (errata row 25)\\
\code{P16-PHIQUANT.md}, \code{P05-CINF.md}, \code{Q15-SPECTRAL.md}, \code{Q04-WEYL.md} & the eight
lemmas: the all-direction floor $\Phi''\succeq(N/12)\mathrm{id}$ and its failure to buy back
$-\beta E[C_{\mathrm{geo}}]$; $E[H_{vv}]\ge0$ for $\theta\le\pi$ with the $\sqrt\beta$ branch beyond;
the Stein identity and the resulting emptiness of the spectral route; the pointwise ratio
$H=\tfrac12\cot(\theta/2)G$ (lemma table of \S\ref{sec:scalar-lemmas})\\
\code{R02-CRIT.md}, \code{R05-WCF.md}, \code{R09-FLUCT.md} & the two moment inequalities with three
independent proofs each, the negative result $E|w|\le\langle w\rangle$ false, the closed forms
$\langle w\rangle_{0}$ and $V_{0}$, the constant $R_{\min}=32.3177$ and the missing monotonicity; also
road 11 of \S\ref{sec:twelve}\\
\code{P21-REBUILD.md}, \code{R07-AFFINE.md} & the skeleton identity and the free-field acid test that
reorganizes the tool set; the two Wronskians, the closed forms for $K_{\mathrm{LS}}$ and
$C_{\mathrm{LS}}$, and the $O(a)$ residual (lemmas 7--8, road 12)\\
\code{P20-THETA.md} & the exact $B$-form $\beta_{0}(B)=B\sqrt{1+9B}$, the product form
$\beta w(\theta)\ge\tfrac19$ and the impossibility of $\theta$-uniformity, the premise collapse at
$\theta=\pi$, the small-$\beta$ leg $\beta\lesssim3.9$, and the small-$B$ bridge broken by its own
counterexample (road 8, \S\ref{sec:window})\\
\code{P19-BETA-AT-L2.md}, \code{Q13-BETA6.md}, \code{P01-L3.md} & the lattice ledger of
\S\ref{sec:lattice-ledger}: the $\beta$ scan $+0.2835/+0.4116/+0.5108/+0.5913$ with $16/16$ chains, the
$\beta=6$ value $+0.6018$ with five estimators and $8/8$ chains, the $\tau_{z}$ correction, and the $L=3$
rerun $+0.3396$ that fixed the fundamental-domain bug\\
\code{R04-L2CERT.md} & the dimension count $192-48-1=143$, the finite-group theorem (orbits are finite,
so symmetry removes no dimension), the $12$-dimensional toron decomposition, the cost of a rigorous
statistical certificate, and the corner law at $a_{0}=16\pi$ (errata row 26)\\
\code{R12-SU3CORR.md}, \code{R11-SU3ADJ.md}, \code{R08-SU3CERT.md} & the $\mathrm{SU}(3)$ verdict
\eqref{eq:su3gap}: the two errors of the first attempt, the adjudication that killed the $S\equiv D$
inference, and the seven-dimensional obstacle\\
\code{Q02-BL-PHI.md}, \code{Q06-RICCATI.md}, \code{Q14-RCBL.md}, \code{Q12-SOS.md},
\code{Q09-COVLOW.md}, \code{Q01-PHI-ONLY.md}, \code{Q08-LARGEA.md}, \code{Q16-CONFLICT.md} & the numbers
of the twelve roads of \S\ref{sec:twelve}, plus the corrected constant $C=3.95$ that supersedes
$1.7454$ and $10.05$\\
\code{docs/paper-v2/main.tex}:30--42 & the original definitions of the label macros
\code{\textbackslash labb}, \code{\textbackslash LT}--\code{\textbackslash LE} (copied here verbatim)\\
\code{CLAUDE.md} \S5 & the current description of the wall (Savvidy concavity, W4's measured $7.1\%$
negative modes)\\
today's $N=3,4$ Monte Carlo & the points $\delta(0),\delta(6),\delta(10)$ of \eqref{eq:n3}, the sign of
$V_0(3)$ and the whole $N$-scaling of \S\ref{sec:failure}: \code{LOOP-RUNBOOK.md} \S11.46--\S11.48, from
\code{research/unconditional/w7/betastar\_n34.log} (\code{betastar\_scaling.py}). Single plaquette/patch
model, not the full theory\\
\end{longtable}
}

\vspace{3pt}
\noindent
\textbf{Deliberately not cited}: any self-published Zenodo / Preprints.org manuscript claiming to have
settled the 4D mass gap (the negative-signal record of \code{MU5-SIGMA-M.md} \S7: they cite each other,
none is peer reviewed, and their ``area law $\Rightarrow$ gap'' step lands on gap L4 unnoticed).
\textbf{What this document did not do}: no new numerics beyond those tabulated in \S\ref{sec:lattice-ledger}
(which are other waves' runs, quoted rather than re-run), no new species, no step declared complete;
progress on Clay is $0\%$. Every verdict above is either ``this class of approach cannot work'' or
``this identity holds'', never ``the gap is proved''. \textbf{One new kind of claim appears in this
edition and only once}: the three enclosed values of \eqref{eq:certpoints}, which are a machine-checked
inequality rather than a measurement --- and they are scoped to one plaquette, one coupling, $a=0$, and
$\mathrm{SU}(2)$.


\end{document}
